\documentclass[11pt]{article}

% ── Core packages ──
\usepackage[utf8]{inputenc}
\usepackage[T1]{fontenc}
\usepackage{lmodern}
\usepackage[margin=1in]{geometry}
\usepackage{amsmath,amssymb}
\usepackage{booktabs}
\usepackage{graphicx}
\usepackage{xcolor}
\usepackage{hyperref}
\usepackage{float}
\usepackage{enumitem}
\usepackage{microtype}
\usepackage{natbib}

% ── Hyperref setup ──
\hypersetup{
  colorlinks=true,
  linkcolor=blue!60!black,
  citecolor=blue!60!black,
  urlcolor=blue!60!black
}

\title{Why PPO Clipping Is Inert in GRPO:\\
Diagnosing the Architectural Disconnection\\
Between Trust Regions and Data Staleness}

\author{
  Yanxin Peng \\
  University of Southern California \\
  \texttt{yanxinpe@usc.edu}
}

\date{\today}

\begin{document}

\maketitle

\begin{abstract}
Group Relative Policy Optimization (GRPO) inherits PPO's clipped surrogate objective as a trust-region mechanism.
Multiple recent works have redesigned this mechanism (DAPO, GSPO, ASPO, CFPO) or removed it entirely (RGRA), citing entropy collapse, training instability, or expendability.
We provide a unifying diagnosis: \textbf{PPO clipping is architecturally inert in standard GRPO training} because the importance ratio uses recomputed log-probabilities, keeping it near unity regardless of data staleness (\texttt{pg\_clipfrac}$\,{=}\,0$).
Using Effective Sample Size (ESS) as a principled offlineness metric, we trace this inactivity to two independent causes: (1)~the PPO loss recomputes its reference log-probs each step, and (2)~conservative learning rates limit per-step policy drift.
We show that reconnecting clipping to data staleness (via stale rollout log-probs) activates the mechanism and produces a $4\times$ entropy increase---but also reveal that activated clipping alone is insufficient for sustained reward improvement over long training horizons.
Our analysis explains why recent clipping redesigns were necessary and why clipping removal was viable, providing a unified lens on an active area of algorithmic development.
\end{abstract}

%% ── Sections ──
%% Sections 1--2: Introduction and Background
%% To be included in main paper via %% Sections 1--2: Introduction and Background
%% To be included in main paper via %% Sections 1--2: Introduction and Background
%% To be included in main paper via \input{sections1_2}

\section{Introduction}
\label{sec:introduction}

Group Relative Policy Optimization \citep[GRPO;][]{deepseekr1} has emerged as the dominant reinforcement learning algorithm for large language model reasoning, powering systems such as DeepSeek-R1 \citep{deepseekr1}, Qwen \citep{qwen2025}, and their derivatives.
GRPO inherits the clipped surrogate objective from Proximal Policy Optimization \citep[PPO;][]{schulman2017ppo}---a trust-region mechanism designed to prevent destructive policy updates by bounding the importance sampling ratio between the current and data-generating policies.
This mechanism is widely regarded as essential to training stability.

Yet a striking pattern has emerged in the recent literature: at least six independent works have redesigned or removed this clipping mechanism, each from a different empirical motivation, and each reporting comparable or improved results.
DAPO \citep{dapo2025} introduces asymmetric clip bounds (\emph{Clip-Higher}) to mitigate entropy collapse.
GSPO \citep{gspo2025} replaces token-level importance sampling with sequence-level ratios, arguing the token-level formulation is ``fundamentally broken.''
CFPO \citep{cfpo2025} substitutes clipping with a quadratic TV-divergence penalty to eliminate dead-zone gradients.
ASPO \citep{aspo2025} inverts the importance ratio for positive-advantage tokens to correct gradient misallocation.
RGRA \citep{rgra2025} removes clipping entirely, retaining plain REINFORCE with group-relative advantages, and matches or exceeds GRPO performance.
\citet{chen2025iclr} analyze clipping as an implicit entropy regularizer and recommend its removal.
All of these modifications work.
The natural question is: \emph{why did each become necessary?}

\paragraph{Our answer.}
We provide a unifying diagnosis.
By instrumenting the PPO loss to track \texttt{pg\_clipfrac}---the fraction of tokens at which the clipped and unclipped objectives differ---we find that this quantity is \emph{exactly zero} across all standard GRPO training configurations we tested.
The clip operator in the PPO surrogate never binds: for every token, the clipped objective equals the unclipped objective, and the trust-region mechanism contributes nothing to the optimization.

The root cause is an architectural disconnection between the importance ratio and data staleness.
In classical PPO, the ratio denominator is the behavior policy that generated the training data.
In GRPO as implemented in standard frameworks (TRL, SLIME, veRL), the denominator uses log-probabilities \emph{recomputed} from the current model at each training step.
Since both numerator and denominator track the same policy, the ratio remains near unity regardless of how far the model has drifted from the policy that generated the rollout data.
We identify two independent causes: (1)~this recomputation of log-probabilities resets the ratio, and (2)~conservative learning rates ($10^{-6}$--$10^{-5}$) limit per-step policy change.
Either cause alone is sufficient to keep clipping inactive; together, they render it entirely inert.

\paragraph{Activating clipping: a controlled experiment.}
We introduce a \texttt{--use-rollout-logprobs} flag that replaces the ratio denominator with the stale rollout log-probabilities, reconnecting the importance ratio to the data-generating distribution.
Combined with tight clipping ($\varepsilon = 0.05$), this is the only configuration in our experimental grid that produces entropy \emph{increase} during training: entropy rises fourfold (from 0.28 to 1.19) while reward initially improves.
However, extended training over 310 rollouts reveals gradual reward decline---the model explores more diverse strategies but does not converge to better ones.
This is itself informative: it demonstrates that clipping is \emph{necessary but not sufficient} under genuine off-policyness, and that complementary mechanisms are needed for sustained improvement.

\paragraph{Contributions.}
Our contributions are:
\begin{enumerate}
    \item The first empirical measurement of clipping inactivity in GRPO training ($\texttt{pg\_clipfrac} = 0$) and identification of its architectural root cause: the PPO importance ratio uses recomputed log-probabilities that are disconnected from data staleness.

    \item The introduction of Effective Sample Size (ESS) as a principled offlineness metric for GRPO, providing a continuous measure of the divergence between the current policy and the data-generating distribution.

    \item A unifying diagnosis that explains why six independent recent works \citep{dapo2025,gspo2025,cfpo2025,aspo2025,rgra2025,chen2025iclr} each found it necessary to redesign or remove GRPO's clipping mechanism.

    \item A controlled factorial experiment showing what happens when clipping is activated: entropy increases and clipping engages as intended, but sustained reward improvement requires mechanisms beyond clipping alone.
\end{enumerate}

\paragraph{Relation to prior observations.}
Several prior works have noted symptoms of clipping inactivity without identifying the cause.
\citet{chen2025iclr} measure $\texttt{pg\_clipfrac} < 0.2\%$ on Qwen2.5-Math-7B at $\text{lr} = 5 \times 10^{-7}$ but attribute the low rate to conservative hyperparameters.
\citet{prosperity2025} report $\texttt{pg\_clipfrac} = 0.05\%$ at staleness $s = 0$ in their multi-policy framework but do not investigate the root cause.
\citet{grpo_secretly_prm} note that the importance ratio equals unity when the number of PPO epochs $\mu = 1$, and \citet{lambert2025} observes that clipping is ``irrelevant'' in this setting.
A TRL user reported $r_t = 1$ in issue \#2769, which the community dismissed as expected behavior.\footnote{\url{https://github.com/huggingface/trl/issues/2769}}
Our contribution is the causal diagnosis---the recomputation of log-probabilities in the PPO loss---and a controlled experiment demonstrating what changes when clipping is made functional.


\section{Background}
\label{sec:background}

\subsection{GRPO and the PPO Clipped Objective}
\label{sec:bg-grpo}

Group Relative Policy Optimization \citep{deepseekr1} generates multiple candidate responses $\{y_1, \dots, y_G\}$ for each prompt $x$ and computes \emph{group-relative advantages}: for each response $y_i$, the advantage $\hat{A}_i$ is derived from the reward $R(x, y_i)$ relative to the group statistics (mean and, optionally, standard deviation) of $\{R(x, y_j)\}_{j=1}^G$.
This eliminates the need for a learned value function, which is a primary source of instability and computational cost in standard PPO for language models.

The policy is then updated using PPO's clipped surrogate objective, applied at the token level.
For a token $a_t$ generated in context $s_t$, the loss is
\begin{equation}
\label{eq:ppo-clip}
\mathcal{L}^{\text{CLIP}}_t(\theta) = -\min\!\Big( r_t(\theta)\,\hat{A}_t,\; \text{clip}\big(r_t(\theta),\, 1{-}\varepsilon,\, 1{+}\varepsilon\big)\,\hat{A}_t \Big),
\end{equation}
where $\varepsilon$ is the clip threshold (typically $0.2$) and $r_t(\theta)$ is the importance sampling ratio
\begin{equation}
\label{eq:is-ratio}
r_t(\theta) = \frac{\pi_\theta(a_t \mid s_t)}{\pi_{\text{old}}(a_t \mid s_t)}.
\end{equation}
When $\hat{A}_t > 0$, the clip operator caps the ratio at $1 + \varepsilon$, preventing excessively large upward updates; when $\hat{A}_t < 0$, it floors the ratio at $1 - \varepsilon$, preventing the policy from moving too far away from dispreferred actions in a single step.
The fraction of tokens at which the clipped and unclipped terms in Equation~\ref{eq:ppo-clip} differ---denoted \texttt{pg\_clipfrac}---measures how often this trust-region constraint is active.

DAPO \citep{dapo2025} modifies this with asymmetric bounds ($\varepsilon_{\text{low}} = 0.2$, $\varepsilon_{\text{high}} = 0.28$), allowing the policy to move further toward preferred actions than away from dispreferred ones.
The standard GRPO training loop additionally includes a per-token KL penalty against a reference policy $\pi_{\text{ref}}$, which we hold fixed in all experiments.


\subsection{The Three Log-Probability References}
\label{sec:bg-three-logprobs}

Understanding why clipping is inert requires distinguishing three log-probability quantities that coexist in GRPO training.
We define them precisely, as conflation of these quantities is the source of the architectural disconnection we diagnose.

\begin{enumerate}
    \item \textbf{Rollout log-probabilities} ($\log \pi_{\text{rollout}}$): computed at data generation time, when the model produces candidate responses for each prompt.
    These represent the actual \emph{behavior policy}---the distribution from which the training data was sampled.
    In code, these are stored as \texttt{batch["rollout\_log\_probs"]}.

    \item \textbf{Old log-probabilities} ($\log \pi_{\text{old}}$): recomputed from the current model at the start of each training step (or PPO epoch).
    In code, these are stored as \texttt{batch["log\_probs"]}.
    In standard GRPO implementations, this quantity serves as the denominator of the importance ratio in Equation~\ref{eq:is-ratio}.

    \item \textbf{Current log-probabilities} ($\log \pi_\theta$): the model's output during the forward pass of the training step.
    This serves as the numerator of the importance ratio.
\end{enumerate}

The critical distinction is the identity of $\pi_{\text{old}}$ in Equation~\ref{eq:is-ratio}.
In classical PPO \citep{schulman2017ppo}, $\pi_{\text{old}} = \pi_{\text{rollout}}$: the ratio denominator is the behavior policy that generated the data, and the ratio measures how far the current policy has moved from the data-generating distribution.
In GRPO as implemented in standard frameworks---including TRL,\footnote{\url{https://github.com/huggingface/trl}} veRL \citep{verl2024}, and SLIME---$\pi_{\text{old}} = \pi_{\theta_{\text{step-start}}}$: the log-probabilities are recomputed from the current model parameters at the beginning of each training step.
This means both the numerator ($\pi_\theta$) and denominator ($\pi_{\text{old}}$) of $r_t(\theta)$ track the same underlying policy, differing only by the parameter updates accumulated during a single gradient step.
The importance ratio is therefore \emph{architecturally near unity}, regardless of how far $\pi_\theta$ has drifted from the rollout policy $\pi_{\text{rollout}}$.

This substitution disconnects the PPO trust region from data staleness.
The clipped objective in Equation~\ref{eq:ppo-clip} constrains how far $\pi_\theta$ moves from $\pi_{\text{old}}$---but since $\pi_{\text{old}}$ is reset to $\pi_\theta$ at each step, the constraint is vacuous.
The policy may drift arbitrarily far from the data-generating distribution $\pi_{\text{rollout}}$ without the clipping mechanism ever detecting or constraining this drift.


\subsection{Offlineness Metrics}
\label{sec:bg-offlineness}

To quantify the divergence between the current policy and the data-generating distribution, we employ the \textbf{Effective Sample Size} (ESS), a standard diagnostic in importance sampling \citep{kong1992ess,liu2001mc}:
\begin{equation}
\label{eq:ess}
\text{ESS} = \frac{\big(\sum_{i} w_i\big)^2}{\sum_{i} w_i^2}, \quad w_i = \frac{\pi_\theta(a_i \mid s_i)}{\pi_{\text{rollout}}(a_i \mid s_i)},
\end{equation}
where the weights $w_i$ are the importance ratios computed against the \emph{rollout} policy (not the recomputed $\pi_{\text{old}}$).
Normalized to the interval $[0, 1]$, $\text{ESS} \approx 1.0$ indicates that the current policy closely matches the data-generating distribution (near on-policy training), while $\text{ESS} \ll 1.0$ indicates severe off-policyness---the training data is unrepresentative of the current policy, and importance-weighted estimates will have high variance.

We additionally track three complementary metrics:
\begin{itemize}
    \item \textbf{Mismatch KL} ($D_{\chi^2}$): a chi-squared divergence proxy measuring distributional mismatch between $\pi_\theta$ and $\pi_{\text{rollout}}$.
    \item \textbf{Clip fraction} (\texttt{clip\_frac\_02}): the fraction of per-token importance weights $w_i$ falling outside the interval $[0.8, 1.2]$, providing a threshold-based staleness indicator.
    \item \textbf{Maximum ratio} (\texttt{max\_ratio}): the largest importance weight $\max_i w_i$, capturing extreme off-policy tokens that may destabilize training.
\end{itemize}

These metrics are computed using rollout log-probabilities as the reference, ensuring they reflect actual data staleness rather than the artificial near-unity ratios produced by the recomputed $\pi_{\text{old}}$.
The distinction between ESS (which uses $\pi_{\text{rollout}}$) and \texttt{pg\_clipfrac} (which uses recomputed $\pi_{\text{old}}$) is central to our diagnosis: a training run can exhibit $\text{ESS} = 0.988$ (measurable off-policyness) while simultaneously showing $\texttt{pg\_clipfrac} = 0$ (no clipping activity), precisely because the two metrics reference different denominators.


\section{Introduction}
\label{sec:introduction}

Group Relative Policy Optimization \citep[GRPO;][]{deepseekr1} has emerged as the dominant reinforcement learning algorithm for large language model reasoning, powering systems such as DeepSeek-R1 \citep{deepseekr1}, Qwen \citep{qwen2025}, and their derivatives.
GRPO inherits the clipped surrogate objective from Proximal Policy Optimization \citep[PPO;][]{schulman2017ppo}---a trust-region mechanism designed to prevent destructive policy updates by bounding the importance sampling ratio between the current and data-generating policies.
This mechanism is widely regarded as essential to training stability.

Yet a striking pattern has emerged in the recent literature: at least six independent works have redesigned or removed this clipping mechanism, each from a different empirical motivation, and each reporting comparable or improved results.
DAPO \citep{dapo2025} introduces asymmetric clip bounds (\emph{Clip-Higher}) to mitigate entropy collapse.
GSPO \citep{gspo2025} replaces token-level importance sampling with sequence-level ratios, arguing the token-level formulation is ``fundamentally broken.''
CFPO \citep{cfpo2025} substitutes clipping with a quadratic TV-divergence penalty to eliminate dead-zone gradients.
ASPO \citep{aspo2025} inverts the importance ratio for positive-advantage tokens to correct gradient misallocation.
RGRA \citep{rgra2025} removes clipping entirely, retaining plain REINFORCE with group-relative advantages, and matches or exceeds GRPO performance.
\citet{chen2025iclr} analyze clipping as an implicit entropy regularizer and recommend its removal.
All of these modifications work.
The natural question is: \emph{why did each become necessary?}

\paragraph{Our answer.}
We provide a unifying diagnosis.
By instrumenting the PPO loss to track \texttt{pg\_clipfrac}---the fraction of tokens at which the clipped and unclipped objectives differ---we find that this quantity is \emph{exactly zero} across all standard GRPO training configurations we tested.
The clip operator in the PPO surrogate never binds: for every token, the clipped objective equals the unclipped objective, and the trust-region mechanism contributes nothing to the optimization.

The root cause is an architectural disconnection between the importance ratio and data staleness.
In classical PPO, the ratio denominator is the behavior policy that generated the training data.
In GRPO as implemented in standard frameworks (TRL, SLIME, veRL), the denominator uses log-probabilities \emph{recomputed} from the current model at each training step.
Since both numerator and denominator track the same policy, the ratio remains near unity regardless of how far the model has drifted from the policy that generated the rollout data.
We identify two independent causes: (1)~this recomputation of log-probabilities resets the ratio, and (2)~conservative learning rates ($10^{-6}$--$10^{-5}$) limit per-step policy change.
Either cause alone is sufficient to keep clipping inactive; together, they render it entirely inert.

\paragraph{Activating clipping: a controlled experiment.}
We introduce a \texttt{--use-rollout-logprobs} flag that replaces the ratio denominator with the stale rollout log-probabilities, reconnecting the importance ratio to the data-generating distribution.
Combined with tight clipping ($\varepsilon = 0.05$), this is the only configuration in our experimental grid that produces entropy \emph{increase} during training: entropy rises fourfold (from 0.28 to 1.19) while reward initially improves.
However, extended training over 310 rollouts reveals gradual reward decline---the model explores more diverse strategies but does not converge to better ones.
This is itself informative: it demonstrates that clipping is \emph{necessary but not sufficient} under genuine off-policyness, and that complementary mechanisms are needed for sustained improvement.

\paragraph{Contributions.}
Our contributions are:
\begin{enumerate}
    \item The first empirical measurement of clipping inactivity in GRPO training ($\texttt{pg\_clipfrac} = 0$) and identification of its architectural root cause: the PPO importance ratio uses recomputed log-probabilities that are disconnected from data staleness.

    \item The introduction of Effective Sample Size (ESS) as a principled offlineness metric for GRPO, providing a continuous measure of the divergence between the current policy and the data-generating distribution.

    \item A unifying diagnosis that explains why six independent recent works \citep{dapo2025,gspo2025,cfpo2025,aspo2025,rgra2025,chen2025iclr} each found it necessary to redesign or remove GRPO's clipping mechanism.

    \item A controlled factorial experiment showing what happens when clipping is activated: entropy increases and clipping engages as intended, but sustained reward improvement requires mechanisms beyond clipping alone.
\end{enumerate}

\paragraph{Relation to prior observations.}
Several prior works have noted symptoms of clipping inactivity without identifying the cause.
\citet{chen2025iclr} measure $\texttt{pg\_clipfrac} < 0.2\%$ on Qwen2.5-Math-7B at $\text{lr} = 5 \times 10^{-7}$ but attribute the low rate to conservative hyperparameters.
\citet{prosperity2025} report $\texttt{pg\_clipfrac} = 0.05\%$ at staleness $s = 0$ in their multi-policy framework but do not investigate the root cause.
\citet{grpo_secretly_prm} note that the importance ratio equals unity when the number of PPO epochs $\mu = 1$, and \citet{lambert2025} observes that clipping is ``irrelevant'' in this setting.
A TRL user reported $r_t = 1$ in issue \#2769, which the community dismissed as expected behavior.\footnote{\url{https://github.com/huggingface/trl/issues/2769}}
Our contribution is the causal diagnosis---the recomputation of log-probabilities in the PPO loss---and a controlled experiment demonstrating what changes when clipping is made functional.


\section{Background}
\label{sec:background}

\subsection{GRPO and the PPO Clipped Objective}
\label{sec:bg-grpo}

Group Relative Policy Optimization \citep{deepseekr1} generates multiple candidate responses $\{y_1, \dots, y_G\}$ for each prompt $x$ and computes \emph{group-relative advantages}: for each response $y_i$, the advantage $\hat{A}_i$ is derived from the reward $R(x, y_i)$ relative to the group statistics (mean and, optionally, standard deviation) of $\{R(x, y_j)\}_{j=1}^G$.
This eliminates the need for a learned value function, which is a primary source of instability and computational cost in standard PPO for language models.

The policy is then updated using PPO's clipped surrogate objective, applied at the token level.
For a token $a_t$ generated in context $s_t$, the loss is
\begin{equation}
\label{eq:ppo-clip}
\mathcal{L}^{\text{CLIP}}_t(\theta) = -\min\!\Big( r_t(\theta)\,\hat{A}_t,\; \text{clip}\big(r_t(\theta),\, 1{-}\varepsilon,\, 1{+}\varepsilon\big)\,\hat{A}_t \Big),
\end{equation}
where $\varepsilon$ is the clip threshold (typically $0.2$) and $r_t(\theta)$ is the importance sampling ratio
\begin{equation}
\label{eq:is-ratio}
r_t(\theta) = \frac{\pi_\theta(a_t \mid s_t)}{\pi_{\text{old}}(a_t \mid s_t)}.
\end{equation}
When $\hat{A}_t > 0$, the clip operator caps the ratio at $1 + \varepsilon$, preventing excessively large upward updates; when $\hat{A}_t < 0$, it floors the ratio at $1 - \varepsilon$, preventing the policy from moving too far away from dispreferred actions in a single step.
The fraction of tokens at which the clipped and unclipped terms in Equation~\ref{eq:ppo-clip} differ---denoted \texttt{pg\_clipfrac}---measures how often this trust-region constraint is active.

DAPO \citep{dapo2025} modifies this with asymmetric bounds ($\varepsilon_{\text{low}} = 0.2$, $\varepsilon_{\text{high}} = 0.28$), allowing the policy to move further toward preferred actions than away from dispreferred ones.
The standard GRPO training loop additionally includes a per-token KL penalty against a reference policy $\pi_{\text{ref}}$, which we hold fixed in all experiments.


\subsection{The Three Log-Probability References}
\label{sec:bg-three-logprobs}

Understanding why clipping is inert requires distinguishing three log-probability quantities that coexist in GRPO training.
We define them precisely, as conflation of these quantities is the source of the architectural disconnection we diagnose.

\begin{enumerate}
    \item \textbf{Rollout log-probabilities} ($\log \pi_{\text{rollout}}$): computed at data generation time, when the model produces candidate responses for each prompt.
    These represent the actual \emph{behavior policy}---the distribution from which the training data was sampled.
    In code, these are stored as \texttt{batch["rollout\_log\_probs"]}.

    \item \textbf{Old log-probabilities} ($\log \pi_{\text{old}}$): recomputed from the current model at the start of each training step (or PPO epoch).
    In code, these are stored as \texttt{batch["log\_probs"]}.
    In standard GRPO implementations, this quantity serves as the denominator of the importance ratio in Equation~\ref{eq:is-ratio}.

    \item \textbf{Current log-probabilities} ($\log \pi_\theta$): the model's output during the forward pass of the training step.
    This serves as the numerator of the importance ratio.
\end{enumerate}

The critical distinction is the identity of $\pi_{\text{old}}$ in Equation~\ref{eq:is-ratio}.
In classical PPO \citep{schulman2017ppo}, $\pi_{\text{old}} = \pi_{\text{rollout}}$: the ratio denominator is the behavior policy that generated the data, and the ratio measures how far the current policy has moved from the data-generating distribution.
In GRPO as implemented in standard frameworks---including TRL,\footnote{\url{https://github.com/huggingface/trl}} veRL \citep{verl2024}, and SLIME---$\pi_{\text{old}} = \pi_{\theta_{\text{step-start}}}$: the log-probabilities are recomputed from the current model parameters at the beginning of each training step.
This means both the numerator ($\pi_\theta$) and denominator ($\pi_{\text{old}}$) of $r_t(\theta)$ track the same underlying policy, differing only by the parameter updates accumulated during a single gradient step.
The importance ratio is therefore \emph{architecturally near unity}, regardless of how far $\pi_\theta$ has drifted from the rollout policy $\pi_{\text{rollout}}$.

This substitution disconnects the PPO trust region from data staleness.
The clipped objective in Equation~\ref{eq:ppo-clip} constrains how far $\pi_\theta$ moves from $\pi_{\text{old}}$---but since $\pi_{\text{old}}$ is reset to $\pi_\theta$ at each step, the constraint is vacuous.
The policy may drift arbitrarily far from the data-generating distribution $\pi_{\text{rollout}}$ without the clipping mechanism ever detecting or constraining this drift.


\subsection{Offlineness Metrics}
\label{sec:bg-offlineness}

To quantify the divergence between the current policy and the data-generating distribution, we employ the \textbf{Effective Sample Size} (ESS), a standard diagnostic in importance sampling \citep{kong1992ess,liu2001mc}:
\begin{equation}
\label{eq:ess}
\text{ESS} = \frac{\big(\sum_{i} w_i\big)^2}{\sum_{i} w_i^2}, \quad w_i = \frac{\pi_\theta(a_i \mid s_i)}{\pi_{\text{rollout}}(a_i \mid s_i)},
\end{equation}
where the weights $w_i$ are the importance ratios computed against the \emph{rollout} policy (not the recomputed $\pi_{\text{old}}$).
Normalized to the interval $[0, 1]$, $\text{ESS} \approx 1.0$ indicates that the current policy closely matches the data-generating distribution (near on-policy training), while $\text{ESS} \ll 1.0$ indicates severe off-policyness---the training data is unrepresentative of the current policy, and importance-weighted estimates will have high variance.

We additionally track three complementary metrics:
\begin{itemize}
    \item \textbf{Mismatch KL} ($D_{\chi^2}$): a chi-squared divergence proxy measuring distributional mismatch between $\pi_\theta$ and $\pi_{\text{rollout}}$.
    \item \textbf{Clip fraction} (\texttt{clip\_frac\_02}): the fraction of per-token importance weights $w_i$ falling outside the interval $[0.8, 1.2]$, providing a threshold-based staleness indicator.
    \item \textbf{Maximum ratio} (\texttt{max\_ratio}): the largest importance weight $\max_i w_i$, capturing extreme off-policy tokens that may destabilize training.
\end{itemize}

These metrics are computed using rollout log-probabilities as the reference, ensuring they reflect actual data staleness rather than the artificial near-unity ratios produced by the recomputed $\pi_{\text{old}}$.
The distinction between ESS (which uses $\pi_{\text{rollout}}$) and \texttt{pg\_clipfrac} (which uses recomputed $\pi_{\text{old}}$) is central to our diagnosis: a training run can exhibit $\text{ESS} = 0.988$ (measurable off-policyness) while simultaneously showing $\texttt{pg\_clipfrac} = 0$ (no clipping activity), precisely because the two metrics reference different denominators.


\section{Introduction}
\label{sec:introduction}

Group Relative Policy Optimization \citep[GRPO;][]{deepseekr1} has emerged as the dominant reinforcement learning algorithm for large language model reasoning, powering systems such as DeepSeek-R1 \citep{deepseekr1}, Qwen \citep{qwen2025}, and their derivatives.
GRPO inherits the clipped surrogate objective from Proximal Policy Optimization \citep[PPO;][]{schulman2017ppo}---a trust-region mechanism designed to prevent destructive policy updates by bounding the importance sampling ratio between the current and data-generating policies.
This mechanism is widely regarded as essential to training stability.

Yet a striking pattern has emerged in the recent literature: at least six independent works have redesigned or removed this clipping mechanism, each from a different empirical motivation, and each reporting comparable or improved results.
DAPO \citep{dapo2025} introduces asymmetric clip bounds (\emph{Clip-Higher}) to mitigate entropy collapse.
GSPO \citep{gspo2025} replaces token-level importance sampling with sequence-level ratios, arguing the token-level formulation is ``fundamentally broken.''
CFPO \citep{cfpo2025} substitutes clipping with a quadratic TV-divergence penalty to eliminate dead-zone gradients.
ASPO \citep{aspo2025} inverts the importance ratio for positive-advantage tokens to correct gradient misallocation.
RGRA \citep{rgra2025} removes clipping entirely, retaining plain REINFORCE with group-relative advantages, and matches or exceeds GRPO performance.
\citet{chen2025iclr} analyze clipping as an implicit entropy regularizer and recommend its removal.
All of these modifications work.
The natural question is: \emph{why did each become necessary?}

\paragraph{Our answer.}
We provide a unifying diagnosis.
By instrumenting the PPO loss to track \texttt{pg\_clipfrac}---the fraction of tokens at which the clipped and unclipped objectives differ---we find that this quantity is \emph{exactly zero} across all standard GRPO training configurations we tested.
The clip operator in the PPO surrogate never binds: for every token, the clipped objective equals the unclipped objective, and the trust-region mechanism contributes nothing to the optimization.

The root cause is an architectural disconnection between the importance ratio and data staleness.
In classical PPO, the ratio denominator is the behavior policy that generated the training data.
In GRPO as implemented in standard frameworks (TRL, SLIME, veRL), the denominator uses log-probabilities \emph{recomputed} from the current model at each training step.
Since both numerator and denominator track the same policy, the ratio remains near unity regardless of how far the model has drifted from the policy that generated the rollout data.
We identify two independent causes: (1)~this recomputation of log-probabilities resets the ratio, and (2)~conservative learning rates ($10^{-6}$--$10^{-5}$) limit per-step policy change.
Either cause alone is sufficient to keep clipping inactive; together, they render it entirely inert.

\paragraph{Activating clipping: a controlled experiment.}
We introduce a \texttt{--use-rollout-logprobs} flag that replaces the ratio denominator with the stale rollout log-probabilities, reconnecting the importance ratio to the data-generating distribution.
Combined with tight clipping ($\varepsilon = 0.05$), this is the only configuration in our experimental grid that produces entropy \emph{increase} during training: entropy rises fourfold (from 0.28 to 1.19) while reward initially improves.
However, extended training over 310 rollouts reveals gradual reward decline---the model explores more diverse strategies but does not converge to better ones.
This is itself informative: it demonstrates that clipping is \emph{necessary but not sufficient} under genuine off-policyness, and that complementary mechanisms are needed for sustained improvement.

\paragraph{Contributions.}
Our contributions are:
\begin{enumerate}
    \item The first empirical measurement of clipping inactivity in GRPO training ($\texttt{pg\_clipfrac} = 0$) and identification of its architectural root cause: the PPO importance ratio uses recomputed log-probabilities that are disconnected from data staleness.

    \item The introduction of Effective Sample Size (ESS) as a principled offlineness metric for GRPO, providing a continuous measure of the divergence between the current policy and the data-generating distribution.

    \item A unifying diagnosis that explains why six independent recent works \citep{dapo2025,gspo2025,cfpo2025,aspo2025,rgra2025,chen2025iclr} each found it necessary to redesign or remove GRPO's clipping mechanism.

    \item A controlled factorial experiment showing what happens when clipping is activated: entropy increases and clipping engages as intended, but sustained reward improvement requires mechanisms beyond clipping alone.
\end{enumerate}

\paragraph{Relation to prior observations.}
Several prior works have noted symptoms of clipping inactivity without identifying the cause.
\citet{chen2025iclr} measure $\texttt{pg\_clipfrac} < 0.2\%$ on Qwen2.5-Math-7B at $\text{lr} = 5 \times 10^{-7}$ but attribute the low rate to conservative hyperparameters.
\citet{prosperity2025} report $\texttt{pg\_clipfrac} = 0.05\%$ at staleness $s = 0$ in their multi-policy framework but do not investigate the root cause.
\citet{grpo_secretly_prm} note that the importance ratio equals unity when the number of PPO epochs $\mu = 1$, and \citet{lambert2025} observes that clipping is ``irrelevant'' in this setting.
A TRL user reported $r_t = 1$ in issue \#2769, which the community dismissed as expected behavior.\footnote{\url{https://github.com/huggingface/trl/issues/2769}}
Our contribution is the causal diagnosis---the recomputation of log-probabilities in the PPO loss---and a controlled experiment demonstrating what changes when clipping is made functional.


\section{Background}
\label{sec:background}

\subsection{GRPO and the PPO Clipped Objective}
\label{sec:bg-grpo}

Group Relative Policy Optimization \citep{deepseekr1} generates multiple candidate responses $\{y_1, \dots, y_G\}$ for each prompt $x$ and computes \emph{group-relative advantages}: for each response $y_i$, the advantage $\hat{A}_i$ is derived from the reward $R(x, y_i)$ relative to the group statistics (mean and, optionally, standard deviation) of $\{R(x, y_j)\}_{j=1}^G$.
This eliminates the need for a learned value function, which is a primary source of instability and computational cost in standard PPO for language models.

The policy is then updated using PPO's clipped surrogate objective, applied at the token level.
For a token $a_t$ generated in context $s_t$, the loss is
\begin{equation}
\label{eq:ppo-clip}
\mathcal{L}^{\text{CLIP}}_t(\theta) = -\min\!\Big( r_t(\theta)\,\hat{A}_t,\; \text{clip}\big(r_t(\theta),\, 1{-}\varepsilon,\, 1{+}\varepsilon\big)\,\hat{A}_t \Big),
\end{equation}
where $\varepsilon$ is the clip threshold (typically $0.2$) and $r_t(\theta)$ is the importance sampling ratio
\begin{equation}
\label{eq:is-ratio}
r_t(\theta) = \frac{\pi_\theta(a_t \mid s_t)}{\pi_{\text{old}}(a_t \mid s_t)}.
\end{equation}
When $\hat{A}_t > 0$, the clip operator caps the ratio at $1 + \varepsilon$, preventing excessively large upward updates; when $\hat{A}_t < 0$, it floors the ratio at $1 - \varepsilon$, preventing the policy from moving too far away from dispreferred actions in a single step.
The fraction of tokens at which the clipped and unclipped terms in Equation~\ref{eq:ppo-clip} differ---denoted \texttt{pg\_clipfrac}---measures how often this trust-region constraint is active.

DAPO \citep{dapo2025} modifies this with asymmetric bounds ($\varepsilon_{\text{low}} = 0.2$, $\varepsilon_{\text{high}} = 0.28$), allowing the policy to move further toward preferred actions than away from dispreferred ones.
The standard GRPO training loop additionally includes a per-token KL penalty against a reference policy $\pi_{\text{ref}}$, which we hold fixed in all experiments.


\subsection{The Three Log-Probability References}
\label{sec:bg-three-logprobs}

Understanding why clipping is inert requires distinguishing three log-probability quantities that coexist in GRPO training.
We define them precisely, as conflation of these quantities is the source of the architectural disconnection we diagnose.

\begin{enumerate}
    \item \textbf{Rollout log-probabilities} ($\log \pi_{\text{rollout}}$): computed at data generation time, when the model produces candidate responses for each prompt.
    These represent the actual \emph{behavior policy}---the distribution from which the training data was sampled.
    In code, these are stored as \texttt{batch["rollout\_log\_probs"]}.

    \item \textbf{Old log-probabilities} ($\log \pi_{\text{old}}$): recomputed from the current model at the start of each training step (or PPO epoch).
    In code, these are stored as \texttt{batch["log\_probs"]}.
    In standard GRPO implementations, this quantity serves as the denominator of the importance ratio in Equation~\ref{eq:is-ratio}.

    \item \textbf{Current log-probabilities} ($\log \pi_\theta$): the model's output during the forward pass of the training step.
    This serves as the numerator of the importance ratio.
\end{enumerate}

The critical distinction is the identity of $\pi_{\text{old}}$ in Equation~\ref{eq:is-ratio}.
In classical PPO \citep{schulman2017ppo}, $\pi_{\text{old}} = \pi_{\text{rollout}}$: the ratio denominator is the behavior policy that generated the data, and the ratio measures how far the current policy has moved from the data-generating distribution.
In GRPO as implemented in standard frameworks---including TRL,\footnote{\url{https://github.com/huggingface/trl}} veRL \citep{verl2024}, and SLIME---$\pi_{\text{old}} = \pi_{\theta_{\text{step-start}}}$: the log-probabilities are recomputed from the current model parameters at the beginning of each training step.
This means both the numerator ($\pi_\theta$) and denominator ($\pi_{\text{old}}$) of $r_t(\theta)$ track the same underlying policy, differing only by the parameter updates accumulated during a single gradient step.
The importance ratio is therefore \emph{architecturally near unity}, regardless of how far $\pi_\theta$ has drifted from the rollout policy $\pi_{\text{rollout}}$.

This substitution disconnects the PPO trust region from data staleness.
The clipped objective in Equation~\ref{eq:ppo-clip} constrains how far $\pi_\theta$ moves from $\pi_{\text{old}}$---but since $\pi_{\text{old}}$ is reset to $\pi_\theta$ at each step, the constraint is vacuous.
The policy may drift arbitrarily far from the data-generating distribution $\pi_{\text{rollout}}$ without the clipping mechanism ever detecting or constraining this drift.


\subsection{Offlineness Metrics}
\label{sec:bg-offlineness}

To quantify the divergence between the current policy and the data-generating distribution, we employ the \textbf{Effective Sample Size} (ESS), a standard diagnostic in importance sampling \citep{kong1992ess,liu2001mc}:
\begin{equation}
\label{eq:ess}
\text{ESS} = \frac{\big(\sum_{i} w_i\big)^2}{\sum_{i} w_i^2}, \quad w_i = \frac{\pi_\theta(a_i \mid s_i)}{\pi_{\text{rollout}}(a_i \mid s_i)},
\end{equation}
where the weights $w_i$ are the importance ratios computed against the \emph{rollout} policy (not the recomputed $\pi_{\text{old}}$).
Normalized to the interval $[0, 1]$, $\text{ESS} \approx 1.0$ indicates that the current policy closely matches the data-generating distribution (near on-policy training), while $\text{ESS} \ll 1.0$ indicates severe off-policyness---the training data is unrepresentative of the current policy, and importance-weighted estimates will have high variance.

We additionally track three complementary metrics:
\begin{itemize}
    \item \textbf{Mismatch KL} ($D_{\chi^2}$): a chi-squared divergence proxy measuring distributional mismatch between $\pi_\theta$ and $\pi_{\text{rollout}}$.
    \item \textbf{Clip fraction} (\texttt{clip\_frac\_02}): the fraction of per-token importance weights $w_i$ falling outside the interval $[0.8, 1.2]$, providing a threshold-based staleness indicator.
    \item \textbf{Maximum ratio} (\texttt{max\_ratio}): the largest importance weight $\max_i w_i$, capturing extreme off-policy tokens that may destabilize training.
\end{itemize}

These metrics are computed using rollout log-probabilities as the reference, ensuring they reflect actual data staleness rather than the artificial near-unity ratios produced by the recomputed $\pi_{\text{old}}$.
The distinction between ESS (which uses $\pi_{\text{rollout}}$) and \texttt{pg\_clipfrac} (which uses recomputed $\pi_{\text{old}}$) is central to our diagnosis: a training run can exhibit $\text{ESS} = 0.988$ (measurable off-policyness) while simultaneously showing $\texttt{pg\_clipfrac} = 0$ (no clipping activity), precisely because the two metrics reference different denominators.

%% Section 3: The Inert Clipping Problem
%% To be included in main paper via %% Section 3: The Inert Clipping Problem
%% To be included in main paper via %% Section 3: The Inert Clipping Problem
%% To be included in main paper via \input{section3}

\section{The Inert Clipping Problem}
\label{sec:inert-clipping}

PPO's clipped surrogate objective is the central trust-region mechanism inherited by GRPO.
In this section, we demonstrate that this mechanism is empirically inert across all standard training configurations we tested, identify two independent root causes, and show how this diagnosis unifies a series of recent algorithmic redesigns.

\subsection{Measuring Clipping Activity}
\label{sec:measuring-clipping}

We instrument the PPO loss to track \texttt{pg\_clipfrac}, the fraction of tokens at which the clipped and unclipped objectives differ---i.e., the fraction of tokens where clipping actively modifies the gradient.
A value of $\texttt{pg\_clipfrac} = 0$ means the clip operator in Equation~\ref{eq:ppo-clip} is never binding: for every token, $\text{clip}(r_t, 1-\varepsilon, 1+\varepsilon) \cdot \hat{A}_t = r_t \cdot \hat{A}_t$, and the trust-region constraint contributes nothing to the optimization.

\paragraph{Setup.}
We train Qwen2.5-Math-1.5B on the DAPO-Math-17K dataset under standard GRPO configurations, resuming from a checkpoint at iteration~799.
All runs use the default clip threshold $\varepsilon = 0.2$.
We additionally track the Effective Sample Size (ESS), defined as
\begin{equation}
\label{eq:ess}
\text{ESS} = \frac{1}{\sum_{i} w_i^2}, \quad w_i = \frac{r_i}{\sum_j r_j},
\end{equation}
where $r_i = \pi_\theta(a_i \mid s) / \pi_{\text{old}}(a_i \mid s)$ are the per-token importance ratios.
ESS measures how many effective independent samples remain after importance weighting; $\text{ESS} \approx 1.0$ indicates near-perfect on-policy training.

\paragraph{PPO epochs do not activate clipping.}
At $\text{lr} = 10^{-6}$, we vary the number of PPO epochs $K \in \{1, 3, 5\}$.
In standard GRPO, multiple PPO epochs recompute log-probabilities at the start of each epoch, resetting the importance ratio denominator.
Across all three settings, $\texttt{pg\_clipfrac} = 0.0$ exactly, and $\text{ESS} = 0.999$.
Increasing the number of optimization passes over the same rollout data does not induce any clipping activity.

\paragraph{Data reuse does not activate clipping.}
We increase the learning rate to $10^{-5}$ and set \texttt{--num-steps-per-rollout}$\,{=}\,8$, which partitions each rollout into 8 non-overlapping training chunks processed sequentially \emph{without} recomputing log-probabilities between steps.
This is the most aggressive data-reuse configuration in our experimental grid.
Even here, $\texttt{pg\_clipfrac} = 0.0$ at every step position, and ESS remains above $0.999$.

\paragraph{Comparison with prior observations.}
Our finding of exactly zero clipping is consistent with, but more extreme than, prior reports.
\citet{chen2025iclr} measure $\texttt{pg\_clipfrac} < 0.2\%$ on Qwen2.5-Math-7B at $\text{lr} = 5 \times 10^{-7}$, attributing the low rate to conservative hyperparameters.
\citet{prosperity2025} report $\texttt{pg\_clipfrac} = 0.05\%$ at staleness $s = 0$ in their multi-policy framework.
Neither work investigates why the rate is near zero or identifies the architectural cause.
The small residual clipping they observe likely reflects numerical differences in their pipelines (larger models, different checkpoint states, or floating-point accumulation effects), but the qualitative conclusion is the same: \emph{clipping is essentially inactive in standard GRPO training}.


\subsection{Root Cause: Two Independent Mechanisms}
\label{sec:root-cause}

We identify two independent mechanisms that jointly render PPO clipping inert.
Their independence is important: removing either one alone is sufficient to produce nonzero clipping activity.

\subsubsection{Cause 1: Architectural Disconnection}
\label{sec:cause-architectural}

The PPO clipped surrogate objective computes the importance ratio as
\begin{equation}
\label{eq:importance-ratio}
r_t(\theta) = \frac{\pi_\theta(a_t \mid s_t)}{\pi_{\text{old}}(a_t \mid s_t)},
\end{equation}
where, in classical PPO, $\pi_{\text{old}}$ represents the \emph{behavior policy} that generated the data.
In GRPO as implemented in standard training frameworks, the denominator uses \texttt{batch["log\_probs"]}---log-probabilities \emph{recomputed} from the current model at the start of each training step---rather than \texttt{batch["rollout\_log\_probs"]}, the log-probabilities recorded at rollout generation time.

This architectural choice means that both numerator and denominator of $r_t(\theta)$ track the same policy (or very nearly so, differing only by the gradient updates within a single step).
Consequently, $r_t(\theta) \approx 1$ regardless of how far the current policy $\pi_\theta$ has drifted from the policy $\pi_{\text{rollout}}$ that actually generated the training data.
The importance ratio is \emph{disconnected} from data staleness: the PPO loss cannot detect, and therefore cannot constrain, policy drift relative to the data-generating distribution.

\subsubsection{Cause 2: Conservative Learning Rates}
\label{sec:cause-lr}

Standard GRPO training uses learning rates in the range $10^{-6}$ to $10^{-5}$.
At these rates, the per-step change in log-probabilities is small enough that even a correctly computed importance ratio (using stale rollout log-probs) would rarely exceed the clip threshold $1 + \varepsilon = 1.2$.
This provides a second, independent barrier to clipping activation: the policy simply does not move fast enough within a single rollout's training steps for the trust region to bind.

\subsubsection{Evidence for Independence}
\label{sec:evidence-independence}

We isolate the two causes by introducing a \texttt{--use-rollout-logprobs} flag that replaces the denominator of $r_t(\theta)$ with the stale rollout log-probabilities, removing Cause~1 while leaving Cause~2 intact.
At $\text{lr} = 10^{-5}$ with \texttt{--num-steps-per-rollout}$\,{=}\,8$:

\begin{itemize}
    \item \textbf{With stale log-probs} (Cause~1 removed): We observe a clear ESS gradient within each rollout, declining from $0.999$ at step~1 to $0.988$ at step~8.
    The variable \texttt{pg\_clipfrac} is nonzero at \emph{all} step positions, rising from $0.001$ (step~1) to $0.024$ (step~8) as accumulated policy drift pushes more tokens past the clip threshold.
    Over 218 rollouts, \texttt{pg\_clipfrac} is consistently positive.

    \item \textbf{Without stale log-probs} (Cause~1 present): Under identical hyperparameters, $\texttt{pg\_clipfrac} = 0.0$ at every step position regardless of $N$.
    The recomputed denominator masks all policy drift.
\end{itemize}

This contrast establishes that Cause~1 and Cause~2 are independent.
Removing the architectural disconnection (Cause~1) is \emph{sufficient} to activate clipping even at the conservative learning rate of $10^{-5}$---the rate merely modulates the \emph{magnitude} of clipping activity, not its presence.
Conversely, increasing the learning rate alone (without fixing the log-prob reference) cannot activate clipping, because the recomputed denominator absorbs all policy changes.


\subsection{Implications for Recent Algorithmic Redesigns}
\label{sec:implications}

Our diagnosis provides a unifying explanation for a striking pattern in recent work: at least six independent papers have redesigned or removed GRPO's clipping mechanism, each motivated by different empirical observations, yet all responding to the same underlying pathology.
Table~\ref{tab:related-redesigns} summarizes how each algorithm's design choice relates to our two-cause diagnosis.

\begin{table*}[t]
\centering
\caption{Recent clipping redesigns interpreted through the lens of inert clipping.
Each modification addresses a symptom of the same root cause: the PPO importance ratio is architecturally near unity, rendering clipping inactive.
``Fresh log-probs'' indicates the algorithm retains the recomputed denominator (Cause~1 unresolved).}
\label{tab:related-redesigns}
\small
\begin{tabular}{@{}p{2.0cm}p{4.5cm}p{7.5cm}@{}}
\toprule
\textbf{Algorithm} & \textbf{Modification} & \textbf{Our Explanation} \\
\midrule
DAPO \newline \citep{dapo2025}
& Asymmetric clip bounds ($\varepsilon_{\text{low}} = 0.2$, $\varepsilon_{\text{high}} = 0.28$)
& Presupposes active clipping.
  If $\texttt{pg\_clipfrac} = 0$, Clip-Higher has no tokens to asymmetrically clip.
  Their reported gains (+8 AIME) are attributable to Dynamic Sampling, not the clip modification (+2 AIME). \\
\addlinespace
GSPO \newline \citep{gspo2025}
& Sequence-level importance sampling replacing token-level
& Retains fresh log-probs; Cause~1 persists.
  Their finding that GSPO clips $100\times$ more tokens than standard GRPO directly supports our diagnosis: GRPO clips almost nothing. \\
\addlinespace
CFPO \newline \citep{cfpo2025}
& Quadratic TV-divergence penalty replacing clip
& Their penalty term is also approximately zero when $r_t \approx 1$.
  The zero-gradient problem they address only manifests after ratios deviate from unity---which requires resolving our upstream Cause~1 first. \\
\addlinespace
ASPO \newline \citep{aspo2025}
& Inverted IS ratios for positive-advantage tokens
& They report average IS weight $\approx 1.0004$ but do not investigate this observation.
  Their gradient-direction fix has negligible magnitude when $r_t \approx 1$; the effective update is dominated by the advantage term alone. \\
\addlinespace
RGRA \newline \citep{rgra2025}
& Remove clipping entirely; use plain REINFORCE with group advantages
& Our finding provides the mechanistic explanation for their empirical result: clipping is removable because it was already absent.
  Removing a no-op preserves performance by construction. \\
\addlinespace
\citet{chen2025iclr}
& Analyze clipping as entropy regularizer; recommend removal
& Their conclusion is correct in the fresh-log-prob regime (Cause~1 present).
  However, under stale log-probs where clipping becomes active, we find it is essential: our no-clip ablation with stale log-probs collapses within 50 rollouts (Section~\ref{sec:factorial}). \\
\bottomrule
\end{tabular}
\end{table*}

The convergent pattern across these works is notable.
Each paper identifies a problem with GRPO's clipping---it does not prevent entropy collapse (DAPO), it operates at the wrong granularity (GSPO), it creates dead zones (CFPO), it misallocates gradients (ASPO), or it is simply unnecessary (RGRA)---and proposes an independent fix.
Our diagnosis reveals these as facets of a single phenomenon: the importance ratio is architecturally constrained to near unity, so the clip operator is never binding.
Under this lens, the modifications fall into two categories: those that \emph{work around} inert clipping by introducing alternative regularization (GSPO, CFPO), and those that \emph{acknowledge} its inertness by simplifying the objective (RGRA, \citealt{chen2025iclr}).
Neither category addresses the upstream cause---the disconnection between the importance ratio and data staleness.

This suggests that a more principled path forward is to \emph{reconnect} the importance ratio to the data-generating policy, making clipping functional as originally intended.
We pursue this direction in Section~\ref{sec:activating-clipping}, where we show that the combination of stale rollout log-probabilities and tight clipping produces a qualitatively different training regime---one in which clipping actively promotes exploration rather than serving as a dormant safety constraint.


\section{The Inert Clipping Problem}
\label{sec:inert-clipping}

PPO's clipped surrogate objective is the central trust-region mechanism inherited by GRPO.
In this section, we demonstrate that this mechanism is empirically inert across all standard training configurations we tested, identify two independent root causes, and show how this diagnosis unifies a series of recent algorithmic redesigns.

\subsection{Measuring Clipping Activity}
\label{sec:measuring-clipping}

We instrument the PPO loss to track \texttt{pg\_clipfrac}, the fraction of tokens at which the clipped and unclipped objectives differ---i.e., the fraction of tokens where clipping actively modifies the gradient.
A value of $\texttt{pg\_clipfrac} = 0$ means the clip operator in Equation~\ref{eq:ppo-clip} is never binding: for every token, $\text{clip}(r_t, 1-\varepsilon, 1+\varepsilon) \cdot \hat{A}_t = r_t \cdot \hat{A}_t$, and the trust-region constraint contributes nothing to the optimization.

\paragraph{Setup.}
We train Qwen2.5-Math-1.5B on the DAPO-Math-17K dataset under standard GRPO configurations, resuming from a checkpoint at iteration~799.
All runs use the default clip threshold $\varepsilon = 0.2$.
We additionally track the Effective Sample Size (ESS), defined as
\begin{equation}
\label{eq:ess}
\text{ESS} = \frac{1}{\sum_{i} w_i^2}, \quad w_i = \frac{r_i}{\sum_j r_j},
\end{equation}
where $r_i = \pi_\theta(a_i \mid s) / \pi_{\text{old}}(a_i \mid s)$ are the per-token importance ratios.
ESS measures how many effective independent samples remain after importance weighting; $\text{ESS} \approx 1.0$ indicates near-perfect on-policy training.

\paragraph{PPO epochs do not activate clipping.}
At $\text{lr} = 10^{-6}$, we vary the number of PPO epochs $K \in \{1, 3, 5\}$.
In standard GRPO, multiple PPO epochs recompute log-probabilities at the start of each epoch, resetting the importance ratio denominator.
Across all three settings, $\texttt{pg\_clipfrac} = 0.0$ exactly, and $\text{ESS} = 0.999$.
Increasing the number of optimization passes over the same rollout data does not induce any clipping activity.

\paragraph{Data reuse does not activate clipping.}
We increase the learning rate to $10^{-5}$ and set \texttt{--num-steps-per-rollout}$\,{=}\,8$, which partitions each rollout into 8 non-overlapping training chunks processed sequentially \emph{without} recomputing log-probabilities between steps.
This is the most aggressive data-reuse configuration in our experimental grid.
Even here, $\texttt{pg\_clipfrac} = 0.0$ at every step position, and ESS remains above $0.999$.

\paragraph{Comparison with prior observations.}
Our finding of exactly zero clipping is consistent with, but more extreme than, prior reports.
\citet{chen2025iclr} measure $\texttt{pg\_clipfrac} < 0.2\%$ on Qwen2.5-Math-7B at $\text{lr} = 5 \times 10^{-7}$, attributing the low rate to conservative hyperparameters.
\citet{prosperity2025} report $\texttt{pg\_clipfrac} = 0.05\%$ at staleness $s = 0$ in their multi-policy framework.
Neither work investigates why the rate is near zero or identifies the architectural cause.
The small residual clipping they observe likely reflects numerical differences in their pipelines (larger models, different checkpoint states, or floating-point accumulation effects), but the qualitative conclusion is the same: \emph{clipping is essentially inactive in standard GRPO training}.


\subsection{Root Cause: Two Independent Mechanisms}
\label{sec:root-cause}

We identify two independent mechanisms that jointly render PPO clipping inert.
Their independence is important: removing either one alone is sufficient to produce nonzero clipping activity.

\subsubsection{Cause 1: Architectural Disconnection}
\label{sec:cause-architectural}

The PPO clipped surrogate objective computes the importance ratio as
\begin{equation}
\label{eq:importance-ratio}
r_t(\theta) = \frac{\pi_\theta(a_t \mid s_t)}{\pi_{\text{old}}(a_t \mid s_t)},
\end{equation}
where, in classical PPO, $\pi_{\text{old}}$ represents the \emph{behavior policy} that generated the data.
In GRPO as implemented in standard training frameworks, the denominator uses \texttt{batch["log\_probs"]}---log-probabilities \emph{recomputed} from the current model at the start of each training step---rather than \texttt{batch["rollout\_log\_probs"]}, the log-probabilities recorded at rollout generation time.

This architectural choice means that both numerator and denominator of $r_t(\theta)$ track the same policy (or very nearly so, differing only by the gradient updates within a single step).
Consequently, $r_t(\theta) \approx 1$ regardless of how far the current policy $\pi_\theta$ has drifted from the policy $\pi_{\text{rollout}}$ that actually generated the training data.
The importance ratio is \emph{disconnected} from data staleness: the PPO loss cannot detect, and therefore cannot constrain, policy drift relative to the data-generating distribution.

\subsubsection{Cause 2: Conservative Learning Rates}
\label{sec:cause-lr}

Standard GRPO training uses learning rates in the range $10^{-6}$ to $10^{-5}$.
At these rates, the per-step change in log-probabilities is small enough that even a correctly computed importance ratio (using stale rollout log-probs) would rarely exceed the clip threshold $1 + \varepsilon = 1.2$.
This provides a second, independent barrier to clipping activation: the policy simply does not move fast enough within a single rollout's training steps for the trust region to bind.

\subsubsection{Evidence for Independence}
\label{sec:evidence-independence}

We isolate the two causes by introducing a \texttt{--use-rollout-logprobs} flag that replaces the denominator of $r_t(\theta)$ with the stale rollout log-probabilities, removing Cause~1 while leaving Cause~2 intact.
At $\text{lr} = 10^{-5}$ with \texttt{--num-steps-per-rollout}$\,{=}\,8$:

\begin{itemize}
    \item \textbf{With stale log-probs} (Cause~1 removed): We observe a clear ESS gradient within each rollout, declining from $0.999$ at step~1 to $0.988$ at step~8.
    The variable \texttt{pg\_clipfrac} is nonzero at \emph{all} step positions, rising from $0.001$ (step~1) to $0.024$ (step~8) as accumulated policy drift pushes more tokens past the clip threshold.
    Over 218 rollouts, \texttt{pg\_clipfrac} is consistently positive.

    \item \textbf{Without stale log-probs} (Cause~1 present): Under identical hyperparameters, $\texttt{pg\_clipfrac} = 0.0$ at every step position regardless of $N$.
    The recomputed denominator masks all policy drift.
\end{itemize}

This contrast establishes that Cause~1 and Cause~2 are independent.
Removing the architectural disconnection (Cause~1) is \emph{sufficient} to activate clipping even at the conservative learning rate of $10^{-5}$---the rate merely modulates the \emph{magnitude} of clipping activity, not its presence.
Conversely, increasing the learning rate alone (without fixing the log-prob reference) cannot activate clipping, because the recomputed denominator absorbs all policy changes.


\subsection{Implications for Recent Algorithmic Redesigns}
\label{sec:implications}

Our diagnosis provides a unifying explanation for a striking pattern in recent work: at least six independent papers have redesigned or removed GRPO's clipping mechanism, each motivated by different empirical observations, yet all responding to the same underlying pathology.
Table~\ref{tab:related-redesigns} summarizes how each algorithm's design choice relates to our two-cause diagnosis.

\begin{table*}[t]
\centering
\caption{Recent clipping redesigns interpreted through the lens of inert clipping.
Each modification addresses a symptom of the same root cause: the PPO importance ratio is architecturally near unity, rendering clipping inactive.
``Fresh log-probs'' indicates the algorithm retains the recomputed denominator (Cause~1 unresolved).}
\label{tab:related-redesigns}
\small
\begin{tabular}{@{}p{2.0cm}p{4.5cm}p{7.5cm}@{}}
\toprule
\textbf{Algorithm} & \textbf{Modification} & \textbf{Our Explanation} \\
\midrule
DAPO \newline \citep{dapo2025}
& Asymmetric clip bounds ($\varepsilon_{\text{low}} = 0.2$, $\varepsilon_{\text{high}} = 0.28$)
& Presupposes active clipping.
  If $\texttt{pg\_clipfrac} = 0$, Clip-Higher has no tokens to asymmetrically clip.
  Their reported gains (+8 AIME) are attributable to Dynamic Sampling, not the clip modification (+2 AIME). \\
\addlinespace
GSPO \newline \citep{gspo2025}
& Sequence-level importance sampling replacing token-level
& Retains fresh log-probs; Cause~1 persists.
  Their finding that GSPO clips $100\times$ more tokens than standard GRPO directly supports our diagnosis: GRPO clips almost nothing. \\
\addlinespace
CFPO \newline \citep{cfpo2025}
& Quadratic TV-divergence penalty replacing clip
& Their penalty term is also approximately zero when $r_t \approx 1$.
  The zero-gradient problem they address only manifests after ratios deviate from unity---which requires resolving our upstream Cause~1 first. \\
\addlinespace
ASPO \newline \citep{aspo2025}
& Inverted IS ratios for positive-advantage tokens
& They report average IS weight $\approx 1.0004$ but do not investigate this observation.
  Their gradient-direction fix has negligible magnitude when $r_t \approx 1$; the effective update is dominated by the advantage term alone. \\
\addlinespace
RGRA \newline \citep{rgra2025}
& Remove clipping entirely; use plain REINFORCE with group advantages
& Our finding provides the mechanistic explanation for their empirical result: clipping is removable because it was already absent.
  Removing a no-op preserves performance by construction. \\
\addlinespace
\citet{chen2025iclr}
& Analyze clipping as entropy regularizer; recommend removal
& Their conclusion is correct in the fresh-log-prob regime (Cause~1 present).
  However, under stale log-probs where clipping becomes active, we find it is essential: our no-clip ablation with stale log-probs collapses within 50 rollouts (Section~\ref{sec:factorial}). \\
\bottomrule
\end{tabular}
\end{table*}

The convergent pattern across these works is notable.
Each paper identifies a problem with GRPO's clipping---it does not prevent entropy collapse (DAPO), it operates at the wrong granularity (GSPO), it creates dead zones (CFPO), it misallocates gradients (ASPO), or it is simply unnecessary (RGRA)---and proposes an independent fix.
Our diagnosis reveals these as facets of a single phenomenon: the importance ratio is architecturally constrained to near unity, so the clip operator is never binding.
Under this lens, the modifications fall into two categories: those that \emph{work around} inert clipping by introducing alternative regularization (GSPO, CFPO), and those that \emph{acknowledge} its inertness by simplifying the objective (RGRA, \citealt{chen2025iclr}).
Neither category addresses the upstream cause---the disconnection between the importance ratio and data staleness.

This suggests that a more principled path forward is to \emph{reconnect} the importance ratio to the data-generating policy, making clipping functional as originally intended.
We pursue this direction in Section~\ref{sec:activating-clipping}, where we show that the combination of stale rollout log-probabilities and tight clipping produces a qualitatively different training regime---one in which clipping actively promotes exploration rather than serving as a dormant safety constraint.


\section{The Inert Clipping Problem}
\label{sec:inert-clipping}

PPO's clipped surrogate objective is the central trust-region mechanism inherited by GRPO.
In this section, we demonstrate that this mechanism is empirically inert across all standard training configurations we tested, identify two independent root causes, and show how this diagnosis unifies a series of recent algorithmic redesigns.

\subsection{Measuring Clipping Activity}
\label{sec:measuring-clipping}

We instrument the PPO loss to track \texttt{pg\_clipfrac}, the fraction of tokens at which the clipped and unclipped objectives differ---i.e., the fraction of tokens where clipping actively modifies the gradient.
A value of $\texttt{pg\_clipfrac} = 0$ means the clip operator in Equation~\ref{eq:ppo-clip} is never binding: for every token, $\text{clip}(r_t, 1-\varepsilon, 1+\varepsilon) \cdot \hat{A}_t = r_t \cdot \hat{A}_t$, and the trust-region constraint contributes nothing to the optimization.

\paragraph{Setup.}
We train Qwen2.5-Math-1.5B on the DAPO-Math-17K dataset under standard GRPO configurations, resuming from a checkpoint at iteration~799.
All runs use the default clip threshold $\varepsilon = 0.2$.
We additionally track the Effective Sample Size (ESS), defined as
\begin{equation}
\label{eq:ess}
\text{ESS} = \frac{1}{\sum_{i} w_i^2}, \quad w_i = \frac{r_i}{\sum_j r_j},
\end{equation}
where $r_i = \pi_\theta(a_i \mid s) / \pi_{\text{old}}(a_i \mid s)$ are the per-token importance ratios.
ESS measures how many effective independent samples remain after importance weighting; $\text{ESS} \approx 1.0$ indicates near-perfect on-policy training.

\paragraph{PPO epochs do not activate clipping.}
At $\text{lr} = 10^{-6}$, we vary the number of PPO epochs $K \in \{1, 3, 5\}$.
In standard GRPO, multiple PPO epochs recompute log-probabilities at the start of each epoch, resetting the importance ratio denominator.
Across all three settings, $\texttt{pg\_clipfrac} = 0.0$ exactly, and $\text{ESS} = 0.999$.
Increasing the number of optimization passes over the same rollout data does not induce any clipping activity.

\paragraph{Data reuse does not activate clipping.}
We increase the learning rate to $10^{-5}$ and set \texttt{--num-steps-per-rollout}$\,{=}\,8$, which partitions each rollout into 8 non-overlapping training chunks processed sequentially \emph{without} recomputing log-probabilities between steps.
This is the most aggressive data-reuse configuration in our experimental grid.
Even here, $\texttt{pg\_clipfrac} = 0.0$ at every step position, and ESS remains above $0.999$.

\paragraph{Comparison with prior observations.}
Our finding of exactly zero clipping is consistent with, but more extreme than, prior reports.
\citet{chen2025iclr} measure $\texttt{pg\_clipfrac} < 0.2\%$ on Qwen2.5-Math-7B at $\text{lr} = 5 \times 10^{-7}$, attributing the low rate to conservative hyperparameters.
\citet{prosperity2025} report $\texttt{pg\_clipfrac} = 0.05\%$ at staleness $s = 0$ in their multi-policy framework.
Neither work investigates why the rate is near zero or identifies the architectural cause.
The small residual clipping they observe likely reflects numerical differences in their pipelines (larger models, different checkpoint states, or floating-point accumulation effects), but the qualitative conclusion is the same: \emph{clipping is essentially inactive in standard GRPO training}.


\subsection{Root Cause: Two Independent Mechanisms}
\label{sec:root-cause}

We identify two independent mechanisms that jointly render PPO clipping inert.
Their independence is important: removing either one alone is sufficient to produce nonzero clipping activity.

\subsubsection{Cause 1: Architectural Disconnection}
\label{sec:cause-architectural}

The PPO clipped surrogate objective computes the importance ratio as
\begin{equation}
\label{eq:importance-ratio}
r_t(\theta) = \frac{\pi_\theta(a_t \mid s_t)}{\pi_{\text{old}}(a_t \mid s_t)},
\end{equation}
where, in classical PPO, $\pi_{\text{old}}$ represents the \emph{behavior policy} that generated the data.
In GRPO as implemented in standard training frameworks, the denominator uses \texttt{batch["log\_probs"]}---log-probabilities \emph{recomputed} from the current model at the start of each training step---rather than \texttt{batch["rollout\_log\_probs"]}, the log-probabilities recorded at rollout generation time.

This architectural choice means that both numerator and denominator of $r_t(\theta)$ track the same policy (or very nearly so, differing only by the gradient updates within a single step).
Consequently, $r_t(\theta) \approx 1$ regardless of how far the current policy $\pi_\theta$ has drifted from the policy $\pi_{\text{rollout}}$ that actually generated the training data.
The importance ratio is \emph{disconnected} from data staleness: the PPO loss cannot detect, and therefore cannot constrain, policy drift relative to the data-generating distribution.

\subsubsection{Cause 2: Conservative Learning Rates}
\label{sec:cause-lr}

Standard GRPO training uses learning rates in the range $10^{-6}$ to $10^{-5}$.
At these rates, the per-step change in log-probabilities is small enough that even a correctly computed importance ratio (using stale rollout log-probs) would rarely exceed the clip threshold $1 + \varepsilon = 1.2$.
This provides a second, independent barrier to clipping activation: the policy simply does not move fast enough within a single rollout's training steps for the trust region to bind.

\subsubsection{Evidence for Independence}
\label{sec:evidence-independence}

We isolate the two causes by introducing a \texttt{--use-rollout-logprobs} flag that replaces the denominator of $r_t(\theta)$ with the stale rollout log-probabilities, removing Cause~1 while leaving Cause~2 intact.
At $\text{lr} = 10^{-5}$ with \texttt{--num-steps-per-rollout}$\,{=}\,8$:

\begin{itemize}
    \item \textbf{With stale log-probs} (Cause~1 removed): We observe a clear ESS gradient within each rollout, declining from $0.999$ at step~1 to $0.988$ at step~8.
    The variable \texttt{pg\_clipfrac} is nonzero at \emph{all} step positions, rising from $0.001$ (step~1) to $0.024$ (step~8) as accumulated policy drift pushes more tokens past the clip threshold.
    Over 218 rollouts, \texttt{pg\_clipfrac} is consistently positive.

    \item \textbf{Without stale log-probs} (Cause~1 present): Under identical hyperparameters, $\texttt{pg\_clipfrac} = 0.0$ at every step position regardless of $N$.
    The recomputed denominator masks all policy drift.
\end{itemize}

This contrast establishes that Cause~1 and Cause~2 are independent.
Removing the architectural disconnection (Cause~1) is \emph{sufficient} to activate clipping even at the conservative learning rate of $10^{-5}$---the rate merely modulates the \emph{magnitude} of clipping activity, not its presence.
Conversely, increasing the learning rate alone (without fixing the log-prob reference) cannot activate clipping, because the recomputed denominator absorbs all policy changes.


\subsection{Implications for Recent Algorithmic Redesigns}
\label{sec:implications}

Our diagnosis provides a unifying explanation for a striking pattern in recent work: at least six independent papers have redesigned or removed GRPO's clipping mechanism, each motivated by different empirical observations, yet all responding to the same underlying pathology.
Table~\ref{tab:related-redesigns} summarizes how each algorithm's design choice relates to our two-cause diagnosis.

\begin{table*}[t]
\centering
\caption{Recent clipping redesigns interpreted through the lens of inert clipping.
Each modification addresses a symptom of the same root cause: the PPO importance ratio is architecturally near unity, rendering clipping inactive.
``Fresh log-probs'' indicates the algorithm retains the recomputed denominator (Cause~1 unresolved).}
\label{tab:related-redesigns}
\small
\begin{tabular}{@{}p{2.0cm}p{4.5cm}p{7.5cm}@{}}
\toprule
\textbf{Algorithm} & \textbf{Modification} & \textbf{Our Explanation} \\
\midrule
DAPO \newline \citep{dapo2025}
& Asymmetric clip bounds ($\varepsilon_{\text{low}} = 0.2$, $\varepsilon_{\text{high}} = 0.28$)
& Presupposes active clipping.
  If $\texttt{pg\_clipfrac} = 0$, Clip-Higher has no tokens to asymmetrically clip.
  Their reported gains (+8 AIME) are attributable to Dynamic Sampling, not the clip modification (+2 AIME). \\
\addlinespace
GSPO \newline \citep{gspo2025}
& Sequence-level importance sampling replacing token-level
& Retains fresh log-probs; Cause~1 persists.
  Their finding that GSPO clips $100\times$ more tokens than standard GRPO directly supports our diagnosis: GRPO clips almost nothing. \\
\addlinespace
CFPO \newline \citep{cfpo2025}
& Quadratic TV-divergence penalty replacing clip
& Their penalty term is also approximately zero when $r_t \approx 1$.
  The zero-gradient problem they address only manifests after ratios deviate from unity---which requires resolving our upstream Cause~1 first. \\
\addlinespace
ASPO \newline \citep{aspo2025}
& Inverted IS ratios for positive-advantage tokens
& They report average IS weight $\approx 1.0004$ but do not investigate this observation.
  Their gradient-direction fix has negligible magnitude when $r_t \approx 1$; the effective update is dominated by the advantage term alone. \\
\addlinespace
RGRA \newline \citep{rgra2025}
& Remove clipping entirely; use plain REINFORCE with group advantages
& Our finding provides the mechanistic explanation for their empirical result: clipping is removable because it was already absent.
  Removing a no-op preserves performance by construction. \\
\addlinespace
\citet{chen2025iclr}
& Analyze clipping as entropy regularizer; recommend removal
& Their conclusion is correct in the fresh-log-prob regime (Cause~1 present).
  However, under stale log-probs where clipping becomes active, we find it is essential: our no-clip ablation with stale log-probs collapses within 50 rollouts (Section~\ref{sec:factorial}). \\
\bottomrule
\end{tabular}
\end{table*}

The convergent pattern across these works is notable.
Each paper identifies a problem with GRPO's clipping---it does not prevent entropy collapse (DAPO), it operates at the wrong granularity (GSPO), it creates dead zones (CFPO), it misallocates gradients (ASPO), or it is simply unnecessary (RGRA)---and proposes an independent fix.
Our diagnosis reveals these as facets of a single phenomenon: the importance ratio is architecturally constrained to near unity, so the clip operator is never binding.
Under this lens, the modifications fall into two categories: those that \emph{work around} inert clipping by introducing alternative regularization (GSPO, CFPO), and those that \emph{acknowledge} its inertness by simplifying the objective (RGRA, \citealt{chen2025iclr}).
Neither category addresses the upstream cause---the disconnection between the importance ratio and data staleness.

This suggests that a more principled path forward is to \emph{reconnect} the importance ratio to the data-generating policy, making clipping functional as originally intended.
We pursue this direction in Section~\ref{sec:activating-clipping}, where we show that the combination of stale rollout log-probabilities and tight clipping produces a qualitatively different training regime---one in which clipping actively promotes exploration rather than serving as a dormant safety constraint.

%% Sections 4, 5, and 6
%% To be included in main paper via %% Sections 4, 5, and 6
%% To be included in main paper via %% Sections 4, 5, and 6
%% To be included in main paper via \input{sections4_5_6}

\section{Activating Clipping: A Controlled Experiment}
\label{sec:activating-clipping}

The diagnosis in Section~\ref{sec:root-cause} suggests a straightforward intervention: replace the recomputed log-probability denominator with the stale rollout log-probabilities, reconnecting the importance ratio to the actual behavior policy.
In this section, we implement this fix, evaluate it through a factorial design, and report both its short-horizon promise and its long-horizon limitations with full transparency.

\subsection{Using Stale Rollout Log-Probabilities}
\label{sec:stale-logprobs}

We introduce a \texttt{--use-rollout-logprobs} flag that modifies the PPO loss computation.
Under the standard implementation, the importance ratio denominator uses \texttt{batch["log\_probs"]}---log-probabilities recomputed from the current policy at each training step.
Our modification substitutes \texttt{batch["rollout\_log\_probs"]}, the log-probabilities recorded at rollout generation time, so that
\begin{equation}
\label{eq:stale-ratio}
r_t(\theta) = \frac{\pi_\theta(a_t \mid s_t)}{\pi_{\text{rollout}}(a_t \mid s_t)},
\end{equation}
where $\pi_{\text{rollout}}$ is the policy that actually generated the training data.
This is a one-line change in the loss computation, but it fundamentally alters the semantics of the importance ratio: $r_t(\theta)$ now measures the full divergence between the current policy and the data-generating policy, rather than the negligible divergence accumulated within a single optimization step.

To prevent the reactivated importance ratio from driving unconstrained policy updates, we pair this modification with a tight clip threshold $\varepsilon = 0.05$, substantially more restrictive than the standard $\varepsilon = 0.2$.
The rationale is that once the importance ratio faithfully reflects policy drift, even modest drift can produce large ratio deviations, and the trust region must tighten accordingly.

\subsection{The 2$\times$2 Factorial Design}
\label{sec:factorial}

To disentangle the contributions of stale log-probabilities and tight clipping, we conduct a $2 \times 2$ factorial experiment crossing two factors: log-probability reference $\in \{\text{fresh}, \text{stale}\}$ and clip threshold $\in \{\varepsilon = 0.05, \varepsilon = 10^6 \text{ (effectively no clip)}\}$.
All conditions use $\text{lr} = 10^{-5}$ with \texttt{--num-steps-per-rollout}$\,{=}\,8$ on the same Qwen2.5-Math-1.5B checkpoint.
Table~\ref{tab:factorial} reports key metrics.

\begin{table*}[t]
\centering
\caption{The $2 \times 2$ factorial: log-probability reference (fresh vs.\ stale) crossed with clip threshold (tight vs.\ none).
ESS and entropy are reported as ranges over the training run (early $\to$ late).
``Collapse'' indicates training instability within the first 50 rollouts.
Only the stale + tight-clip condition produces increased exploration without collapse.}
\label{tab:factorial}
\small
\begin{tabular}{@{}llcccc@{}}
\toprule
\textbf{Log-probs} & \textbf{Clip $\varepsilon$} & \textbf{ESS} & \textbf{Entropy} & \textbf{Reward} & \textbf{pg\_clipfrac} \\
\midrule
Fresh & $0.05$ & $0.999$ & $0.20 \to 0.24$ & $0.13 \to 0.15$ & $0.0$ \\
Fresh & $10^6$ (none) & \multicolumn{3}{c}{Collapse within $\sim$50 rollouts} & N/A \\
Stale & $0.05$ & $0.999 \to 0.997$ & $0.28 \to 1.19$ ($4\times$) & $0.18 \to 0.20$ & $0.01 \to 0.04$ \\
Stale & $10^6$ (none) & $0.956$--$0.988$ & Unstable & Degrading & N/A \\
\bottomrule
\end{tabular}
\end{table*}

Several findings emerge from this design:

\paragraph{Synergistic interaction, not additive effects.}
Neither factor alone produces the entropy increase observed in the stale + tight-clip condition.
Fresh log-probabilities with tight clipping (row~1) yield flat entropy ($0.20 \to 0.24$) and zero clipping activity---the tight threshold cannot bind when the importance ratio is architecturally near unity.
Stale log-probabilities without clipping (row~4) produce an unstable training trajectory with degrading reward, as unconstrained importance ratios amplify off-policy errors.
Only when both factors are combined does a qualitatively different regime emerge: entropy increases fourfold ($0.28 \to 1.19$) while ESS remains above $0.997$ and reward improves.
This is a genuine interaction effect, not an additive combination.

\paragraph{Clipping is active at all step positions.}
In the stale + tight-clip condition, \texttt{pg\_clipfrac} is nonzero at every step position within each rollout, rising from $0.001$ at step~1 to $0.024$ at step~8.
This gradient across step positions reflects the progressive accumulation of policy drift within a rollout: later training steps operate on data that is increasingly stale relative to the current policy, activating the clip constraint more frequently.
Over 186 rollouts, clipping activity rises from $1.4\%$ to $4.4\%$ as training progresses, confirming that the trust region is genuinely binding and modulating the optimization.

\paragraph{Clipping is necessary under off-policyness.}
The fresh + no-clip condition collapses within approximately 50 rollouts, while the stale + no-clip condition exhibits ESS values as low as $0.956$ with unstable dynamics.
Clipping---when actually active---provides a meaningful stability guarantee that is absent in the standard (inert) configuration only because the on-policy regime makes it unnecessary.

\subsection{Long-Horizon Behavior}
\label{sec:long-horizon}

The 186-rollout factorial results suggest that activated clipping promotes exploration and improves reward.
To assess whether these gains are sustained, we extend the stale + tight-clip condition to 310 rollouts (approximately 1,100 training steps at $\text{lr} = 10^{-5}$).

\paragraph{Stability is maintained.}
ESS remains in the range $0.988$--$0.999$ throughout the extended run, with no off-policy collapse.
This confirms that tight clipping ($\varepsilon = 0.05$) effectively constrains policy drift even over substantially longer training horizons.

\paragraph{Exploration persists.}
Entropy continues to exhibit high-entropy cycling with spikes reaching $1.37$, indicating sustained exploration of diverse strategies.
The model does not converge to a narrow mode, in contrast to the entropy-decreasing trajectories typical of standard GRPO training.

\paragraph{Reward does not sustain improvement.}
Reward peaks at $0.305$ around rollout~834, then gradually declines to $0.02$--$0.07$ by rollout~1,110.
Concurrently, the truncation rate climbs from $14\%$ to $48\%$, indicating that the model increasingly generates responses that exceed the maximum sequence length without producing valid solutions.

\paragraph{Interpretation.}
Activated clipping successfully promotes exploration and provides an initial reward improvement, confirming that the trust region mechanism is functioning as intended.
However, it is insufficient for sustained learning: the model explores more diverse strategies but does not reliably converge to better ones.
The reward decline suggests that exploration without a complementary convergence mechanism---such as a learned value function, reward shaping, or curriculum design---eventually leads the model into low-reward regions of the policy space from which group-relative advantages alone cannot recover.

We view this as an informative boundary result rather than a failure of the approach.
It delineates what PPO clipping can and cannot accomplish when properly activated: clipping constrains the \emph{rate} of policy change per step, but it does not guide the \emph{direction} of exploration.
In the standard GRPO setting, this limitation is invisible because clipping never fires.
Activating it reveals the need for additional mechanisms that standard on-policy GRPO happens to avoid by not exploring aggressively.

\subsection{Learning Rate Sensitivity}
\label{sec:lr-sensitivity}

The stale-logprobs fix does not broaden GRPO's stability regime with respect to learning rate.
We evaluate two additional configurations at $\text{lr} = 5 \times 10^{-5}$:

\begin{itemize}
    \item \textbf{$\varepsilon = 0.05$:} Entropy explodes within 50 rollouts, and the model degenerates into repetitive outputs (e.g., \texttt{"n n n \ldots"}), with reward collapsing to zero.
    The learning rate induces per-step policy changes so large that even tight clipping cannot prevent mode collapse.

    \item \textbf{$\varepsilon = 0.02$:} Entropy remains controlled, but the model suffers capability degradation, with reward declining from $0.05$ to $0.008$.
    The extremely tight trust region prevents collapse but also prevents the model from making meaningful policy improvements.
\end{itemize}

These results establish that the stale-logprobs fix operates within GRPO's existing learning rate stability regime ($\text{lr} \leq 10^{-5}$), not as a license for more aggressive optimization.
At learning rates where standard GRPO is already unstable, reconnecting the importance ratio amplifies rather than mitigates the instability, because the trust region now faithfully reports the (excessive) policy drift that conservative learning rates were implicitly suppressing.


\section{Discussion}
\label{sec:discussion}

\subsection{Why Does GRPO Work Without Active Clipping?}
\label{sec:why-grpo-works}

If PPO clipping is architecturally inert, what provides GRPO's implicit regularization?
Recent theoretical work offers a compelling answer.
\citet{grpo-dpo2025} show that GRPO's group-relative objective is equivalent to an implicit contrastive loss, closely related to Direct Preference Optimization (DPO).
Under this interpretation, the group-relative advantages---computed by normalizing rewards within each sample group---already implement a form of pairwise preference learning: the model is trained to increase the probability of higher-reward completions relative to lower-reward ones within the same group, rather than to maximize absolute reward.

This contrastive structure provides implicit regularization through two mechanisms.
First, the zero-mean property of group-relative advantages ensures that the net gradient across a group is balanced---some responses are upweighted while others are downweighted---preventing the runaway probability mass accumulation that PPO's trust region was designed to prevent.
Second, the within-group normalization ties the magnitude of policy updates to the reward \emph{variance} within each group, naturally attenuating updates when the model is already performing uniformly (all completions receive similar rewards).

This explains two otherwise puzzling observations.
RGRA~\citep{rgra2025} can remove clipping entirely without performance loss, because the contrastive objective already provides the stability guarantee that clipping was supposed to enforce.
Conversely, our finding that the no-clip ablation collapses under stale log-probabilities (Section~\ref{sec:factorial}) is consistent with this explanation: when the importance ratio is disconnected from data staleness, the contrastive structure is sufficient; when genuinely off-policy data is introduced, it is not, and an explicit trust region becomes necessary.

\subsection{When Does Clipping Matter?}
\label{sec:when-clipping-matters}

Our results identify a precise boundary: clipping becomes necessary when training data is genuinely off-policy relative to the current model.
This condition arises in several practical scenarios:

\begin{itemize}
    \item \textbf{Aggressive data reuse:} Replaying rollout buffers across multiple training iterations, as in experience replay variants of RLHF.
    \item \textbf{Asynchronous training:} When rollout generation and policy optimization proceed concurrently, the policy may update substantially between data generation and training.
    \item \textbf{Multi-policy frameworks:} \citet{prosperity2025} report that at staleness $s = 256$, clipping activity rises to $1.22\%$---still low, but nonzero and increasing with staleness.
\end{itemize}

Our stale-logprobs experiment introduces controlled off-policyness in a single-policy setting, providing a clean test of clipping's role.
The result is unambiguous: without clipping, stale-logprob training collapses within approximately 50 rollouts; with tight clipping, training remains stable for over 300 rollouts.
Clipping is not merely a theoretical safeguard under off-policyness---it is an empirically necessary one.

\subsection{Practical Implications}
\label{sec:practical-implications}

Our findings have distinct implications for two categories of practitioners:

\paragraph{Standard GRPO users.}
If you train with fresh rollouts and conservative learning rates ($\text{lr} \leq 10^{-5}$), your clipping mechanism is doing nothing.
This is acceptable---the contrastive structure of group-relative advantages provides sufficient regularization, and removing clipping (as in RGRA) would produce identical training dynamics.
Practitioners in this regime should not expect clipping-related hyperparameters ($\varepsilon$, asymmetric bounds) to affect training outcomes.

\paragraph{Aggressive data-reuse settings.}
If your training pipeline introduces genuine off-policyness---through stale rollout buffers, asynchronous updates, or multi-policy sampling---you need a trust-region mechanism that actually functions.
Several options exist:
\begin{enumerate}[label=(\alph*)]
    \item \textbf{Stale log-probs with tight clipping} (our approach): simple to implement as a one-line change, but with limited long-term reward improvement (Section~\ref{sec:long-horizon}).
    \item \textbf{Sequence-level importance sampling} (GSPO; \citealt{gspo2025}): operates at the response level rather than token level, avoiding the high variance of token-level IS.
    \item \textbf{Second-moment masking} (M2PO; \citealt{m2po2025}): uses variance of importance ratios to identify and mask unreliable tokens.
    \item \textbf{Quadratic divergence penalties} (CFPO; \citealt{cfpo2025}): replaces the hard clip with a smooth penalty that avoids the zero-gradient problem.
\end{enumerate}
The choice among these approaches depends on the specific source of off-policyness and the acceptable implementation complexity.
Our contribution is diagnostic rather than prescriptive: we establish \emph{why} the standard mechanism fails, enabling informed selection among alternatives.

\subsection{Limitations}
\label{sec:limitations}

Several limitations constrain the generality of our findings.

\paragraph{Single model and domain.}
All experiments use Qwen2.5-Math-1.5B on mathematical reasoning tasks.
While our core finding (zero clipping activity) is consistent with measurements by \citet{chen2025iclr} on Qwen2.5-Math-7B and by \citet{prosperity2025} across multiple model scales and domains, the long-horizon behavior of the stale-logprobs fix has not been validated at larger scales or on non-mathematical tasks.

\paragraph{Reward decline under activated clipping.}
The stale + tight-clip configuration shows reward decline after approximately 800 rollouts (Section~\ref{sec:long-horizon}), limiting its practical value as a standalone fix.
We are transparent that this represents an incomplete solution to the off-policy GRPO problem.

\paragraph{No head-to-head comparison.}
We do not directly compare the stale-logprobs fix against DAPO, GSPO, or other recent algorithms on the same benchmark, making it difficult to rank the approaches in terms of downstream task performance.

\paragraph{Checkpoint artifacts.}
All experiments resume from a pre-trained checkpoint at iteration~799, which carries optimizer state from prior training.
While we control for learning-rate contamination via \texttt{--no-load-optim}, residual effects of the checkpoint's training history cannot be fully excluded.


\section{Conclusion}
\label{sec:conclusion}

PPO clipping in GRPO is architecturally inert.
Across all standard configurations we tested---varying PPO epochs, learning rates, and data-reuse strategies---the clipped surrogate objective reduces to the unclipped objective exactly, with the clip operator never binding.
We trace this inertness to two independent causes: (1)~the recomputation of log-probabilities in the PPO loss, which disconnects the importance ratio from data staleness, and (2)~conservative learning rates that keep per-step policy changes below the clip threshold even when the ratio is correctly computed.

This diagnosis provides a unifying explanation for the proliferation of clipping redesigns in recent work.
DAPO's asymmetric bounds, GSPO's sequence-level importance sampling, CFPO's quadratic penalty, ASPO's inverted ratios, and RGRA's outright removal of clipping are all responses---from different angles---to the same underlying pathology: a trust-region mechanism that never engages.
The redesigns succeed not because they improve clipping, but because they either replace it with an alternative that functions (GSPO, CFPO) or acknowledge its redundancy (RGRA).

We show that reconnecting the importance ratio to the data-generating policy via stale rollout log-probabilities, combined with a tight clip threshold, activates the trust region and produces a qualitatively different training regime---one characterized by sustained exploration and initial reward improvement.
However, activated clipping alone is insufficient for long-term learning: reward declines after approximately 800 rollouts, revealing that the trust region constrains the rate of policy change but does not guide its direction.
This points to the need for complementary mechanisms---value functions, reward shaping, or curriculum design---in any off-policy extension of GRPO.

The broader lesson is methodological.
When a safety mechanism is structurally prevented from firing, its presence creates a false sense of security while its hyperparameters become unidentifiable nuisance parameters.
Diagnosing this structural inertness is a prerequisite for principled algorithm design, whether the goal is to fix the mechanism, replace it, or remove it.


\section{Activating Clipping: A Controlled Experiment}
\label{sec:activating-clipping}

The diagnosis in Section~\ref{sec:root-cause} suggests a straightforward intervention: replace the recomputed log-probability denominator with the stale rollout log-probabilities, reconnecting the importance ratio to the actual behavior policy.
In this section, we implement this fix, evaluate it through a factorial design, and report both its short-horizon promise and its long-horizon limitations with full transparency.

\subsection{Using Stale Rollout Log-Probabilities}
\label{sec:stale-logprobs}

We introduce a \texttt{--use-rollout-logprobs} flag that modifies the PPO loss computation.
Under the standard implementation, the importance ratio denominator uses \texttt{batch["log\_probs"]}---log-probabilities recomputed from the current policy at each training step.
Our modification substitutes \texttt{batch["rollout\_log\_probs"]}, the log-probabilities recorded at rollout generation time, so that
\begin{equation}
\label{eq:stale-ratio}
r_t(\theta) = \frac{\pi_\theta(a_t \mid s_t)}{\pi_{\text{rollout}}(a_t \mid s_t)},
\end{equation}
where $\pi_{\text{rollout}}$ is the policy that actually generated the training data.
This is a one-line change in the loss computation, but it fundamentally alters the semantics of the importance ratio: $r_t(\theta)$ now measures the full divergence between the current policy and the data-generating policy, rather than the negligible divergence accumulated within a single optimization step.

To prevent the reactivated importance ratio from driving unconstrained policy updates, we pair this modification with a tight clip threshold $\varepsilon = 0.05$, substantially more restrictive than the standard $\varepsilon = 0.2$.
The rationale is that once the importance ratio faithfully reflects policy drift, even modest drift can produce large ratio deviations, and the trust region must tighten accordingly.

\subsection{The 2$\times$2 Factorial Design}
\label{sec:factorial}

To disentangle the contributions of stale log-probabilities and tight clipping, we conduct a $2 \times 2$ factorial experiment crossing two factors: log-probability reference $\in \{\text{fresh}, \text{stale}\}$ and clip threshold $\in \{\varepsilon = 0.05, \varepsilon = 10^6 \text{ (effectively no clip)}\}$.
All conditions use $\text{lr} = 10^{-5}$ with \texttt{--num-steps-per-rollout}$\,{=}\,8$ on the same Qwen2.5-Math-1.5B checkpoint.
Table~\ref{tab:factorial} reports key metrics.

\begin{table*}[t]
\centering
\caption{The $2 \times 2$ factorial: log-probability reference (fresh vs.\ stale) crossed with clip threshold (tight vs.\ none).
ESS and entropy are reported as ranges over the training run (early $\to$ late).
``Collapse'' indicates training instability within the first 50 rollouts.
Only the stale + tight-clip condition produces increased exploration without collapse.}
\label{tab:factorial}
\small
\begin{tabular}{@{}llcccc@{}}
\toprule
\textbf{Log-probs} & \textbf{Clip $\varepsilon$} & \textbf{ESS} & \textbf{Entropy} & \textbf{Reward} & \textbf{pg\_clipfrac} \\
\midrule
Fresh & $0.05$ & $0.999$ & $0.20 \to 0.24$ & $0.13 \to 0.15$ & $0.0$ \\
Fresh & $10^6$ (none) & \multicolumn{3}{c}{Collapse within $\sim$50 rollouts} & N/A \\
Stale & $0.05$ & $0.999 \to 0.997$ & $0.28 \to 1.19$ ($4\times$) & $0.18 \to 0.20$ & $0.01 \to 0.04$ \\
Stale & $10^6$ (none) & $0.956$--$0.988$ & Unstable & Degrading & N/A \\
\bottomrule
\end{tabular}
\end{table*}

Several findings emerge from this design:

\paragraph{Synergistic interaction, not additive effects.}
Neither factor alone produces the entropy increase observed in the stale + tight-clip condition.
Fresh log-probabilities with tight clipping (row~1) yield flat entropy ($0.20 \to 0.24$) and zero clipping activity---the tight threshold cannot bind when the importance ratio is architecturally near unity.
Stale log-probabilities without clipping (row~4) produce an unstable training trajectory with degrading reward, as unconstrained importance ratios amplify off-policy errors.
Only when both factors are combined does a qualitatively different regime emerge: entropy increases fourfold ($0.28 \to 1.19$) while ESS remains above $0.997$ and reward improves.
This is a genuine interaction effect, not an additive combination.

\paragraph{Clipping is active at all step positions.}
In the stale + tight-clip condition, \texttt{pg\_clipfrac} is nonzero at every step position within each rollout, rising from $0.001$ at step~1 to $0.024$ at step~8.
This gradient across step positions reflects the progressive accumulation of policy drift within a rollout: later training steps operate on data that is increasingly stale relative to the current policy, activating the clip constraint more frequently.
Over 186 rollouts, clipping activity rises from $1.4\%$ to $4.4\%$ as training progresses, confirming that the trust region is genuinely binding and modulating the optimization.

\paragraph{Clipping is necessary under off-policyness.}
The fresh + no-clip condition collapses within approximately 50 rollouts, while the stale + no-clip condition exhibits ESS values as low as $0.956$ with unstable dynamics.
Clipping---when actually active---provides a meaningful stability guarantee that is absent in the standard (inert) configuration only because the on-policy regime makes it unnecessary.

\subsection{Long-Horizon Behavior}
\label{sec:long-horizon}

The 186-rollout factorial results suggest that activated clipping promotes exploration and improves reward.
To assess whether these gains are sustained, we extend the stale + tight-clip condition to 310 rollouts (approximately 1,100 training steps at $\text{lr} = 10^{-5}$).

\paragraph{Stability is maintained.}
ESS remains in the range $0.988$--$0.999$ throughout the extended run, with no off-policy collapse.
This confirms that tight clipping ($\varepsilon = 0.05$) effectively constrains policy drift even over substantially longer training horizons.

\paragraph{Exploration persists.}
Entropy continues to exhibit high-entropy cycling with spikes reaching $1.37$, indicating sustained exploration of diverse strategies.
The model does not converge to a narrow mode, in contrast to the entropy-decreasing trajectories typical of standard GRPO training.

\paragraph{Reward does not sustain improvement.}
Reward peaks at $0.305$ around rollout~834, then gradually declines to $0.02$--$0.07$ by rollout~1,110.
Concurrently, the truncation rate climbs from $14\%$ to $48\%$, indicating that the model increasingly generates responses that exceed the maximum sequence length without producing valid solutions.

\paragraph{Interpretation.}
Activated clipping successfully promotes exploration and provides an initial reward improvement, confirming that the trust region mechanism is functioning as intended.
However, it is insufficient for sustained learning: the model explores more diverse strategies but does not reliably converge to better ones.
The reward decline suggests that exploration without a complementary convergence mechanism---such as a learned value function, reward shaping, or curriculum design---eventually leads the model into low-reward regions of the policy space from which group-relative advantages alone cannot recover.

We view this as an informative boundary result rather than a failure of the approach.
It delineates what PPO clipping can and cannot accomplish when properly activated: clipping constrains the \emph{rate} of policy change per step, but it does not guide the \emph{direction} of exploration.
In the standard GRPO setting, this limitation is invisible because clipping never fires.
Activating it reveals the need for additional mechanisms that standard on-policy GRPO happens to avoid by not exploring aggressively.

\subsection{Learning Rate Sensitivity}
\label{sec:lr-sensitivity}

The stale-logprobs fix does not broaden GRPO's stability regime with respect to learning rate.
We evaluate two additional configurations at $\text{lr} = 5 \times 10^{-5}$:

\begin{itemize}
    \item \textbf{$\varepsilon = 0.05$:} Entropy explodes within 50 rollouts, and the model degenerates into repetitive outputs (e.g., \texttt{"n n n \ldots"}), with reward collapsing to zero.
    The learning rate induces per-step policy changes so large that even tight clipping cannot prevent mode collapse.

    \item \textbf{$\varepsilon = 0.02$:} Entropy remains controlled, but the model suffers capability degradation, with reward declining from $0.05$ to $0.008$.
    The extremely tight trust region prevents collapse but also prevents the model from making meaningful policy improvements.
\end{itemize}

These results establish that the stale-logprobs fix operates within GRPO's existing learning rate stability regime ($\text{lr} \leq 10^{-5}$), not as a license for more aggressive optimization.
At learning rates where standard GRPO is already unstable, reconnecting the importance ratio amplifies rather than mitigates the instability, because the trust region now faithfully reports the (excessive) policy drift that conservative learning rates were implicitly suppressing.


\section{Discussion}
\label{sec:discussion}

\subsection{Why Does GRPO Work Without Active Clipping?}
\label{sec:why-grpo-works}

If PPO clipping is architecturally inert, what provides GRPO's implicit regularization?
Recent theoretical work offers a compelling answer.
\citet{grpo-dpo2025} show that GRPO's group-relative objective is equivalent to an implicit contrastive loss, closely related to Direct Preference Optimization (DPO).
Under this interpretation, the group-relative advantages---computed by normalizing rewards within each sample group---already implement a form of pairwise preference learning: the model is trained to increase the probability of higher-reward completions relative to lower-reward ones within the same group, rather than to maximize absolute reward.

This contrastive structure provides implicit regularization through two mechanisms.
First, the zero-mean property of group-relative advantages ensures that the net gradient across a group is balanced---some responses are upweighted while others are downweighted---preventing the runaway probability mass accumulation that PPO's trust region was designed to prevent.
Second, the within-group normalization ties the magnitude of policy updates to the reward \emph{variance} within each group, naturally attenuating updates when the model is already performing uniformly (all completions receive similar rewards).

This explains two otherwise puzzling observations.
RGRA~\citep{rgra2025} can remove clipping entirely without performance loss, because the contrastive objective already provides the stability guarantee that clipping was supposed to enforce.
Conversely, our finding that the no-clip ablation collapses under stale log-probabilities (Section~\ref{sec:factorial}) is consistent with this explanation: when the importance ratio is disconnected from data staleness, the contrastive structure is sufficient; when genuinely off-policy data is introduced, it is not, and an explicit trust region becomes necessary.

\subsection{When Does Clipping Matter?}
\label{sec:when-clipping-matters}

Our results identify a precise boundary: clipping becomes necessary when training data is genuinely off-policy relative to the current model.
This condition arises in several practical scenarios:

\begin{itemize}
    \item \textbf{Aggressive data reuse:} Replaying rollout buffers across multiple training iterations, as in experience replay variants of RLHF.
    \item \textbf{Asynchronous training:} When rollout generation and policy optimization proceed concurrently, the policy may update substantially between data generation and training.
    \item \textbf{Multi-policy frameworks:} \citet{prosperity2025} report that at staleness $s = 256$, clipping activity rises to $1.22\%$---still low, but nonzero and increasing with staleness.
\end{itemize}

Our stale-logprobs experiment introduces controlled off-policyness in a single-policy setting, providing a clean test of clipping's role.
The result is unambiguous: without clipping, stale-logprob training collapses within approximately 50 rollouts; with tight clipping, training remains stable for over 300 rollouts.
Clipping is not merely a theoretical safeguard under off-policyness---it is an empirically necessary one.

\subsection{Practical Implications}
\label{sec:practical-implications}

Our findings have distinct implications for two categories of practitioners:

\paragraph{Standard GRPO users.}
If you train with fresh rollouts and conservative learning rates ($\text{lr} \leq 10^{-5}$), your clipping mechanism is doing nothing.
This is acceptable---the contrastive structure of group-relative advantages provides sufficient regularization, and removing clipping (as in RGRA) would produce identical training dynamics.
Practitioners in this regime should not expect clipping-related hyperparameters ($\varepsilon$, asymmetric bounds) to affect training outcomes.

\paragraph{Aggressive data-reuse settings.}
If your training pipeline introduces genuine off-policyness---through stale rollout buffers, asynchronous updates, or multi-policy sampling---you need a trust-region mechanism that actually functions.
Several options exist:
\begin{enumerate}[label=(\alph*)]
    \item \textbf{Stale log-probs with tight clipping} (our approach): simple to implement as a one-line change, but with limited long-term reward improvement (Section~\ref{sec:long-horizon}).
    \item \textbf{Sequence-level importance sampling} (GSPO; \citealt{gspo2025}): operates at the response level rather than token level, avoiding the high variance of token-level IS.
    \item \textbf{Second-moment masking} (M2PO; \citealt{m2po2025}): uses variance of importance ratios to identify and mask unreliable tokens.
    \item \textbf{Quadratic divergence penalties} (CFPO; \citealt{cfpo2025}): replaces the hard clip with a smooth penalty that avoids the zero-gradient problem.
\end{enumerate}
The choice among these approaches depends on the specific source of off-policyness and the acceptable implementation complexity.
Our contribution is diagnostic rather than prescriptive: we establish \emph{why} the standard mechanism fails, enabling informed selection among alternatives.

\subsection{Limitations}
\label{sec:limitations}

Several limitations constrain the generality of our findings.

\paragraph{Single model and domain.}
All experiments use Qwen2.5-Math-1.5B on mathematical reasoning tasks.
While our core finding (zero clipping activity) is consistent with measurements by \citet{chen2025iclr} on Qwen2.5-Math-7B and by \citet{prosperity2025} across multiple model scales and domains, the long-horizon behavior of the stale-logprobs fix has not been validated at larger scales or on non-mathematical tasks.

\paragraph{Reward decline under activated clipping.}
The stale + tight-clip configuration shows reward decline after approximately 800 rollouts (Section~\ref{sec:long-horizon}), limiting its practical value as a standalone fix.
We are transparent that this represents an incomplete solution to the off-policy GRPO problem.

\paragraph{No head-to-head comparison.}
We do not directly compare the stale-logprobs fix against DAPO, GSPO, or other recent algorithms on the same benchmark, making it difficult to rank the approaches in terms of downstream task performance.

\paragraph{Checkpoint artifacts.}
All experiments resume from a pre-trained checkpoint at iteration~799, which carries optimizer state from prior training.
While we control for learning-rate contamination via \texttt{--no-load-optim}, residual effects of the checkpoint's training history cannot be fully excluded.


\section{Conclusion}
\label{sec:conclusion}

PPO clipping in GRPO is architecturally inert.
Across all standard configurations we tested---varying PPO epochs, learning rates, and data-reuse strategies---the clipped surrogate objective reduces to the unclipped objective exactly, with the clip operator never binding.
We trace this inertness to two independent causes: (1)~the recomputation of log-probabilities in the PPO loss, which disconnects the importance ratio from data staleness, and (2)~conservative learning rates that keep per-step policy changes below the clip threshold even when the ratio is correctly computed.

This diagnosis provides a unifying explanation for the proliferation of clipping redesigns in recent work.
DAPO's asymmetric bounds, GSPO's sequence-level importance sampling, CFPO's quadratic penalty, ASPO's inverted ratios, and RGRA's outright removal of clipping are all responses---from different angles---to the same underlying pathology: a trust-region mechanism that never engages.
The redesigns succeed not because they improve clipping, but because they either replace it with an alternative that functions (GSPO, CFPO) or acknowledge its redundancy (RGRA).

We show that reconnecting the importance ratio to the data-generating policy via stale rollout log-probabilities, combined with a tight clip threshold, activates the trust region and produces a qualitatively different training regime---one characterized by sustained exploration and initial reward improvement.
However, activated clipping alone is insufficient for long-term learning: reward declines after approximately 800 rollouts, revealing that the trust region constrains the rate of policy change but does not guide its direction.
This points to the need for complementary mechanisms---value functions, reward shaping, or curriculum design---in any off-policy extension of GRPO.

The broader lesson is methodological.
When a safety mechanism is structurally prevented from firing, its presence creates a false sense of security while its hyperparameters become unidentifiable nuisance parameters.
Diagnosing this structural inertness is a prerequisite for principled algorithm design, whether the goal is to fix the mechanism, replace it, or remove it.


\section{Activating Clipping: A Controlled Experiment}
\label{sec:activating-clipping}

The diagnosis in Section~\ref{sec:root-cause} suggests a straightforward intervention: replace the recomputed log-probability denominator with the stale rollout log-probabilities, reconnecting the importance ratio to the actual behavior policy.
In this section, we implement this fix, evaluate it through a factorial design, and report both its short-horizon promise and its long-horizon limitations with full transparency.

\subsection{Using Stale Rollout Log-Probabilities}
\label{sec:stale-logprobs}

We introduce a \texttt{--use-rollout-logprobs} flag that modifies the PPO loss computation.
Under the standard implementation, the importance ratio denominator uses \texttt{batch["log\_probs"]}---log-probabilities recomputed from the current policy at each training step.
Our modification substitutes \texttt{batch["rollout\_log\_probs"]}, the log-probabilities recorded at rollout generation time, so that
\begin{equation}
\label{eq:stale-ratio}
r_t(\theta) = \frac{\pi_\theta(a_t \mid s_t)}{\pi_{\text{rollout}}(a_t \mid s_t)},
\end{equation}
where $\pi_{\text{rollout}}$ is the policy that actually generated the training data.
This is a one-line change in the loss computation, but it fundamentally alters the semantics of the importance ratio: $r_t(\theta)$ now measures the full divergence between the current policy and the data-generating policy, rather than the negligible divergence accumulated within a single optimization step.

To prevent the reactivated importance ratio from driving unconstrained policy updates, we pair this modification with a tight clip threshold $\varepsilon = 0.05$, substantially more restrictive than the standard $\varepsilon = 0.2$.
The rationale is that once the importance ratio faithfully reflects policy drift, even modest drift can produce large ratio deviations, and the trust region must tighten accordingly.

\subsection{The 2$\times$2 Factorial Design}
\label{sec:factorial}

To disentangle the contributions of stale log-probabilities and tight clipping, we conduct a $2 \times 2$ factorial experiment crossing two factors: log-probability reference $\in \{\text{fresh}, \text{stale}\}$ and clip threshold $\in \{\varepsilon = 0.05, \varepsilon = 10^6 \text{ (effectively no clip)}\}$.
All conditions use $\text{lr} = 10^{-5}$ with \texttt{--num-steps-per-rollout}$\,{=}\,8$ on the same Qwen2.5-Math-1.5B checkpoint.
Table~\ref{tab:factorial} reports key metrics.

\begin{table*}[t]
\centering
\caption{The $2 \times 2$ factorial: log-probability reference (fresh vs.\ stale) crossed with clip threshold (tight vs.\ none).
ESS and entropy are reported as ranges over the training run (early $\to$ late).
``Collapse'' indicates training instability within the first 50 rollouts.
Only the stale + tight-clip condition produces increased exploration without collapse.}
\label{tab:factorial}
\small
\begin{tabular}{@{}llcccc@{}}
\toprule
\textbf{Log-probs} & \textbf{Clip $\varepsilon$} & \textbf{ESS} & \textbf{Entropy} & \textbf{Reward} & \textbf{pg\_clipfrac} \\
\midrule
Fresh & $0.05$ & $0.999$ & $0.20 \to 0.24$ & $0.13 \to 0.15$ & $0.0$ \\
Fresh & $10^6$ (none) & \multicolumn{3}{c}{Collapse within $\sim$50 rollouts} & N/A \\
Stale & $0.05$ & $0.999 \to 0.997$ & $0.28 \to 1.19$ ($4\times$) & $0.18 \to 0.20$ & $0.01 \to 0.04$ \\
Stale & $10^6$ (none) & $0.956$--$0.988$ & Unstable & Degrading & N/A \\
\bottomrule
\end{tabular}
\end{table*}

Several findings emerge from this design:

\paragraph{Synergistic interaction, not additive effects.}
Neither factor alone produces the entropy increase observed in the stale + tight-clip condition.
Fresh log-probabilities with tight clipping (row~1) yield flat entropy ($0.20 \to 0.24$) and zero clipping activity---the tight threshold cannot bind when the importance ratio is architecturally near unity.
Stale log-probabilities without clipping (row~4) produce an unstable training trajectory with degrading reward, as unconstrained importance ratios amplify off-policy errors.
Only when both factors are combined does a qualitatively different regime emerge: entropy increases fourfold ($0.28 \to 1.19$) while ESS remains above $0.997$ and reward improves.
This is a genuine interaction effect, not an additive combination.

\paragraph{Clipping is active at all step positions.}
In the stale + tight-clip condition, \texttt{pg\_clipfrac} is nonzero at every step position within each rollout, rising from $0.001$ at step~1 to $0.024$ at step~8.
This gradient across step positions reflects the progressive accumulation of policy drift within a rollout: later training steps operate on data that is increasingly stale relative to the current policy, activating the clip constraint more frequently.
Over 186 rollouts, clipping activity rises from $1.4\%$ to $4.4\%$ as training progresses, confirming that the trust region is genuinely binding and modulating the optimization.

\paragraph{Clipping is necessary under off-policyness.}
The fresh + no-clip condition collapses within approximately 50 rollouts, while the stale + no-clip condition exhibits ESS values as low as $0.956$ with unstable dynamics.
Clipping---when actually active---provides a meaningful stability guarantee that is absent in the standard (inert) configuration only because the on-policy regime makes it unnecessary.

\subsection{Long-Horizon Behavior}
\label{sec:long-horizon}

The 186-rollout factorial results suggest that activated clipping promotes exploration and improves reward.
To assess whether these gains are sustained, we extend the stale + tight-clip condition to 310 rollouts (approximately 1,100 training steps at $\text{lr} = 10^{-5}$).

\paragraph{Stability is maintained.}
ESS remains in the range $0.988$--$0.999$ throughout the extended run, with no off-policy collapse.
This confirms that tight clipping ($\varepsilon = 0.05$) effectively constrains policy drift even over substantially longer training horizons.

\paragraph{Exploration persists.}
Entropy continues to exhibit high-entropy cycling with spikes reaching $1.37$, indicating sustained exploration of diverse strategies.
The model does not converge to a narrow mode, in contrast to the entropy-decreasing trajectories typical of standard GRPO training.

\paragraph{Reward does not sustain improvement.}
Reward peaks at $0.305$ around rollout~834, then gradually declines to $0.02$--$0.07$ by rollout~1,110.
Concurrently, the truncation rate climbs from $14\%$ to $48\%$, indicating that the model increasingly generates responses that exceed the maximum sequence length without producing valid solutions.

\paragraph{Interpretation.}
Activated clipping successfully promotes exploration and provides an initial reward improvement, confirming that the trust region mechanism is functioning as intended.
However, it is insufficient for sustained learning: the model explores more diverse strategies but does not reliably converge to better ones.
The reward decline suggests that exploration without a complementary convergence mechanism---such as a learned value function, reward shaping, or curriculum design---eventually leads the model into low-reward regions of the policy space from which group-relative advantages alone cannot recover.

We view this as an informative boundary result rather than a failure of the approach.
It delineates what PPO clipping can and cannot accomplish when properly activated: clipping constrains the \emph{rate} of policy change per step, but it does not guide the \emph{direction} of exploration.
In the standard GRPO setting, this limitation is invisible because clipping never fires.
Activating it reveals the need for additional mechanisms that standard on-policy GRPO happens to avoid by not exploring aggressively.

\subsection{Learning Rate Sensitivity}
\label{sec:lr-sensitivity}

The stale-logprobs fix does not broaden GRPO's stability regime with respect to learning rate.
We evaluate two additional configurations at $\text{lr} = 5 \times 10^{-5}$:

\begin{itemize}
    \item \textbf{$\varepsilon = 0.05$:} Entropy explodes within 50 rollouts, and the model degenerates into repetitive outputs (e.g., \texttt{"n n n \ldots"}), with reward collapsing to zero.
    The learning rate induces per-step policy changes so large that even tight clipping cannot prevent mode collapse.

    \item \textbf{$\varepsilon = 0.02$:} Entropy remains controlled, but the model suffers capability degradation, with reward declining from $0.05$ to $0.008$.
    The extremely tight trust region prevents collapse but also prevents the model from making meaningful policy improvements.
\end{itemize}

These results establish that the stale-logprobs fix operates within GRPO's existing learning rate stability regime ($\text{lr} \leq 10^{-5}$), not as a license for more aggressive optimization.
At learning rates where standard GRPO is already unstable, reconnecting the importance ratio amplifies rather than mitigates the instability, because the trust region now faithfully reports the (excessive) policy drift that conservative learning rates were implicitly suppressing.


\section{Discussion}
\label{sec:discussion}

\subsection{Why Does GRPO Work Without Active Clipping?}
\label{sec:why-grpo-works}

If PPO clipping is architecturally inert, what provides GRPO's implicit regularization?
Recent theoretical work offers a compelling answer.
\citet{grpo-dpo2025} show that GRPO's group-relative objective is equivalent to an implicit contrastive loss, closely related to Direct Preference Optimization (DPO).
Under this interpretation, the group-relative advantages---computed by normalizing rewards within each sample group---already implement a form of pairwise preference learning: the model is trained to increase the probability of higher-reward completions relative to lower-reward ones within the same group, rather than to maximize absolute reward.

This contrastive structure provides implicit regularization through two mechanisms.
First, the zero-mean property of group-relative advantages ensures that the net gradient across a group is balanced---some responses are upweighted while others are downweighted---preventing the runaway probability mass accumulation that PPO's trust region was designed to prevent.
Second, the within-group normalization ties the magnitude of policy updates to the reward \emph{variance} within each group, naturally attenuating updates when the model is already performing uniformly (all completions receive similar rewards).

This explains two otherwise puzzling observations.
RGRA~\citep{rgra2025} can remove clipping entirely without performance loss, because the contrastive objective already provides the stability guarantee that clipping was supposed to enforce.
Conversely, our finding that the no-clip ablation collapses under stale log-probabilities (Section~\ref{sec:factorial}) is consistent with this explanation: when the importance ratio is disconnected from data staleness, the contrastive structure is sufficient; when genuinely off-policy data is introduced, it is not, and an explicit trust region becomes necessary.

\subsection{When Does Clipping Matter?}
\label{sec:when-clipping-matters}

Our results identify a precise boundary: clipping becomes necessary when training data is genuinely off-policy relative to the current model.
This condition arises in several practical scenarios:

\begin{itemize}
    \item \textbf{Aggressive data reuse:} Replaying rollout buffers across multiple training iterations, as in experience replay variants of RLHF.
    \item \textbf{Asynchronous training:} When rollout generation and policy optimization proceed concurrently, the policy may update substantially between data generation and training.
    \item \textbf{Multi-policy frameworks:} \citet{prosperity2025} report that at staleness $s = 256$, clipping activity rises to $1.22\%$---still low, but nonzero and increasing with staleness.
\end{itemize}

Our stale-logprobs experiment introduces controlled off-policyness in a single-policy setting, providing a clean test of clipping's role.
The result is unambiguous: without clipping, stale-logprob training collapses within approximately 50 rollouts; with tight clipping, training remains stable for over 300 rollouts.
Clipping is not merely a theoretical safeguard under off-policyness---it is an empirically necessary one.

\subsection{Practical Implications}
\label{sec:practical-implications}

Our findings have distinct implications for two categories of practitioners:

\paragraph{Standard GRPO users.}
If you train with fresh rollouts and conservative learning rates ($\text{lr} \leq 10^{-5}$), your clipping mechanism is doing nothing.
This is acceptable---the contrastive structure of group-relative advantages provides sufficient regularization, and removing clipping (as in RGRA) would produce identical training dynamics.
Practitioners in this regime should not expect clipping-related hyperparameters ($\varepsilon$, asymmetric bounds) to affect training outcomes.

\paragraph{Aggressive data-reuse settings.}
If your training pipeline introduces genuine off-policyness---through stale rollout buffers, asynchronous updates, or multi-policy sampling---you need a trust-region mechanism that actually functions.
Several options exist:
\begin{enumerate}[label=(\alph*)]
    \item \textbf{Stale log-probs with tight clipping} (our approach): simple to implement as a one-line change, but with limited long-term reward improvement (Section~\ref{sec:long-horizon}).
    \item \textbf{Sequence-level importance sampling} (GSPO; \citealt{gspo2025}): operates at the response level rather than token level, avoiding the high variance of token-level IS.
    \item \textbf{Second-moment masking} (M2PO; \citealt{m2po2025}): uses variance of importance ratios to identify and mask unreliable tokens.
    \item \textbf{Quadratic divergence penalties} (CFPO; \citealt{cfpo2025}): replaces the hard clip with a smooth penalty that avoids the zero-gradient problem.
\end{enumerate}
The choice among these approaches depends on the specific source of off-policyness and the acceptable implementation complexity.
Our contribution is diagnostic rather than prescriptive: we establish \emph{why} the standard mechanism fails, enabling informed selection among alternatives.

\subsection{Limitations}
\label{sec:limitations}

Several limitations constrain the generality of our findings.

\paragraph{Single model and domain.}
All experiments use Qwen2.5-Math-1.5B on mathematical reasoning tasks.
While our core finding (zero clipping activity) is consistent with measurements by \citet{chen2025iclr} on Qwen2.5-Math-7B and by \citet{prosperity2025} across multiple model scales and domains, the long-horizon behavior of the stale-logprobs fix has not been validated at larger scales or on non-mathematical tasks.

\paragraph{Reward decline under activated clipping.}
The stale + tight-clip configuration shows reward decline after approximately 800 rollouts (Section~\ref{sec:long-horizon}), limiting its practical value as a standalone fix.
We are transparent that this represents an incomplete solution to the off-policy GRPO problem.

\paragraph{No head-to-head comparison.}
We do not directly compare the stale-logprobs fix against DAPO, GSPO, or other recent algorithms on the same benchmark, making it difficult to rank the approaches in terms of downstream task performance.

\paragraph{Checkpoint artifacts.}
All experiments resume from a pre-trained checkpoint at iteration~799, which carries optimizer state from prior training.
While we control for learning-rate contamination via \texttt{--no-load-optim}, residual effects of the checkpoint's training history cannot be fully excluded.


\section{Conclusion}
\label{sec:conclusion}

PPO clipping in GRPO is architecturally inert.
Across all standard configurations we tested---varying PPO epochs, learning rates, and data-reuse strategies---the clipped surrogate objective reduces to the unclipped objective exactly, with the clip operator never binding.
We trace this inertness to two independent causes: (1)~the recomputation of log-probabilities in the PPO loss, which disconnects the importance ratio from data staleness, and (2)~conservative learning rates that keep per-step policy changes below the clip threshold even when the ratio is correctly computed.

This diagnosis provides a unifying explanation for the proliferation of clipping redesigns in recent work.
DAPO's asymmetric bounds, GSPO's sequence-level importance sampling, CFPO's quadratic penalty, ASPO's inverted ratios, and RGRA's outright removal of clipping are all responses---from different angles---to the same underlying pathology: a trust-region mechanism that never engages.
The redesigns succeed not because they improve clipping, but because they either replace it with an alternative that functions (GSPO, CFPO) or acknowledge its redundancy (RGRA).

We show that reconnecting the importance ratio to the data-generating policy via stale rollout log-probabilities, combined with a tight clip threshold, activates the trust region and produces a qualitatively different training regime---one characterized by sustained exploration and initial reward improvement.
However, activated clipping alone is insufficient for long-term learning: reward declines after approximately 800 rollouts, revealing that the trust region constrains the rate of policy change but does not guide its direction.
This points to the need for complementary mechanisms---value functions, reward shaping, or curriculum design---in any off-policy extension of GRPO.

The broader lesson is methodological.
When a safety mechanism is structurally prevented from firing, its presence creates a false sense of security while its hyperparameters become unidentifiable nuisance parameters.
Diagnosing this structural inertness is a prerequisite for principled algorithm design, whether the goal is to fix the mechanism, replace it, or remove it.


%% ── References ──
% TODO: Create references.bib with proper keys
% \bibliographystyle{plainnat}
% \bibliography{references}

\end{document}
